%!TEX root = groundunitbriefing.tex

% LTeX: enabled=true

% LTeX: dictionary+=aéronefs
% LTeX: dictionary+=AFVs
% LTeX: dictionary+=APCs
% LTeX: dictionary+=Ayn-al-Saqr
% LTeX: dictionary+=Biryusa
% LTeX: dictionary+=BOFI
% LTeX: dictionary+=BOFI-R
% LTeX: dictionary+=Bofors
% LTeX: dictionary+=contre
% LTeX: dictionary+=Dochka
% LTeX: dictionary+=DShK
% LTeX: dictionary+=Dushka
% LTeX: dictionary+=EOTG
% LTeX: dictionary+=Fliegerabwehr
% LTeX: dictionary+=Fuzes
% LTeX: dictionary+=Gemag
% LTeX: dictionary+=Gepard
% LTeX: dictionary+=Hispano-Suiza
% LTeX: dictionary+=Janes
% LTeX: dictionary+=Kanone
% LTeX: dictionary+=Luchs
% LTeX: dictionary+=Malvinas
% LTeX: dictionary+=Marder
% LTeX: dictionary+=Mk
% LTeX: dictionary+=MobileRadar
% LTeX: dictionary+=modèle
% LTeX: dictionary+=Multisensor
% LTeX: dictionary+=Oeil
% LTeX: dictionary+=Oerlikon
% LTeX: dictionary+=Ogaden
% LTeX: dictionary+=Panhard
% LTeX: dictionary+=PIVADS
% LTeX: dictionary+=Proximity-fuzed
% LTeX: dictionary+=RadarTutorial
% LTeX: dictionary+=SAMs
% LTeX: dictionary+=Skyguard
% LTeX: dictionary+=Skysweeper
% LTeX: dictionary+=TSOH
% LTeX: dictionary+=VADS
% LTeX: dictionary+=VDAA
% LTeX: dictionary+=Vz
% LTeX: dictionary+=WeapoNEWS
% LTeX: dictionary+=WeaponSystems
% LTeX: dictionary+=Werrell

% Note: The British spellings "Defense" is toggled locally when needed to allow
% checking the US spelling elsewhere.

\chapter{AAA Units}

AAA units use guns to destroy attacking aircraft, but also have a secondary ground-combat capability. Some also have SAMs, giving a layered capacity. AAA units are presented below with their counters, ordered first by caliber and then by year of introduction.

The \hyperlink{table:aaa-units}{AAA Unit Table} summarizes the properties of AAA units.
As well as the basic properties common to all ground units, the table gives:
the AAA class (light, medium, or heavy);
the maximum altitude;
the limits of short, medium, and long range in hexes;
the hit rolls at short, medium, and long range;
the damage rating;
whether the unit has an all-weather or ranging-only fire-control radar and its frequency band or a laser rangefinder;
whether it has night infrared sights; and
whether it is also equipped with a SAM system.
If the unit has a SAM, this will be described in detail in the section on SAM launcher units.%
\note{AAA Values}{
    This long note summarizes how I have determined the AAA values for AAA units. VP values are discussed separately.

    I compare AAA values for AAA units from {\AirStr}, {\EOTG}, {\TSOH}, and Malcolm Pipe's 2026 \href{https://airpower.groups.io/g/main/files/Puma\%27s\%20Files/Charts\%20&\%20Tables/Pods&Ordnance-2-2026.xlsx}{compilation} to the actual characteristics of the guns they represent. My aim is to empirically understand the underlying model sufficiently to identify anomalous values and to apply it to other guns.

    \noteparagraph{AAA Values for Batteries} I start by attempting to understand the AAA values for batteries.

    The upper part of the following table shows the original data for gun batteries from {\AirStr} (AS), {\EOTG} (EOTG), {\TSOH} (TSOH). I also include data for these guns from Malcolm Pipes' 2026 compilation (MP). For the ZPU-2, ZPU-4, Twin Rh-202, and KS-19, these sources are not in agreement among themselves and so there are multiple entries. I have added columns with the caliber, burst rate of fire per mount (i.e., the rate of fire per barrel multiplied by the number of barrels per mount), and muzzle velocity. Values in red will be subsequently modified.

    \begin{minipage}{\linewidth}
        \setlength{\tabcolsep}{0.2em}
        \begin{center}
            \smallskip
            AAA VALUES FOR BATTERIES\par
            \medskip
            \scriptsize
            \begin{tabularx}{\linewidth}{LrrlrclrrL}
                \toprule
                Type
                &\multicolumn{1}{c}{\vertical{Maximum Altitude}}
                &\multicolumn{1}{c}{\vertical{Ranges}}
                &\multicolumn{1}{c}{\vertical{Hit Rolls}}
                &\multicolumn{1}{c}{\vertical{Damage}}
                &\multicolumn{1}{c}{\vertical{FCR/LRF}}
                &\multicolumn{1}{c}{\vertical{Caliber (mm)}}
                &\multicolumn{1}{c}{\vertical{Rate of Fire (RPM)}}
                &\multicolumn{1}{c}{\vertical{Muzzle Velocity (m/s)}}
                &\multicolumn{1}{c}{\vertical{Sources}}
                \\
                \midrule
                \multicolumn{10}{c}{Original Values}\\
                \midrule
                M2&5&2/3/4&3/2/1&1&---&\phantom{0}12.7&500&890&AS/MP\\
                DShK&5&2/3/4&3/2/1&1&---&\phantom{0}12.7&500&850&AS/MP\\
                M55&5&2/3/4&5/4/3&2&---&\phantom{0}12.7&2000&890&EOTG/MP\\
                ZPU-1&6&2/3/5&3/2/1&1&---&\phantom{0}14.5&600&1000&EOTG/TSOH/MP\\
                ZPU-2&6&2/3/5&4/3/2&1&---&\phantom{0}14.5&1200&1000&AS\\
                ZPU-2&6&2/3/5&\highlight{4/3/1}&\highlight{3}&---&\phantom{0}14.5&1200&1000&MP\\
                ZPU-4&6&2/3/5&5/4/3&2&---&\phantom{0}14.5&2400&1000&AS/MP\\
                ZPU-4&6&2/3/5&\highlight{4/3/2}&2&---&\phantom{0}14.5&2400&1000&TSOH\\
                TCM-20&8&2/4/6&4/4/3&2&---&\phantom{0}20&1450&860&EOTG\\
                Single Rh-202&8&2/4/6&\highlight{3/2/1}&2&---&\phantom{0}20&1030&1044&AS\\
                Twin Rh-202&8&2/4/6&5/4/3&\highlight{2}&---&\phantom{0}20&2060&1044&AS\\
                Twin Rh-202&8&2/4/6&\highlight{3/2/1}&3&---&\phantom{0}20&2060&1044&MP\\
                M167&8&2/4/6&6/5/4&3&R&\phantom{0}20&3000&1030&AS/MP\\
                ZU-23&9&2/5/8&\highlight{4/3/2}&\highlight{2}&---&\phantom{0}23&2000&980&AS/TSOH/MP\\
                Gemag&10&3/6/9&6/5/4&4&R&\phantom{0}25&1800&1040&AS\\
                Oerlikon GDF-001&12&4/8/10&\highlight{5/4/3}&4&---&\phantom{0}35&1100&1175&AS/MP\\
                61-K (M-38)&12&4/8/11&3/2/1&4&---&\phantom{0}37&165&880&AS/TSOH/MP\\
                Bofors L/60&15&4/8/12&2/2/1&4&---&\phantom{0}40&135&865&EOTG/MP\\
                Bofors L/70&15&4/8/12&\highlight{3/3/2}&4&---&\phantom{0}40&330&1000&AS/MP\\
                BOFI&15&4/8/12&\highlight{4/4/3}&4&W&\phantom{0}40&330&1000&AS\\
                S-60&18&5/10/15&2/2/1&5&---&\phantom{0}57&110&1000&AS/TSOH/MP\\
                KS-12&27&6/12/18&\highlight{2/1/0}&6&---&\phantom{0}85&10&790&AS/TSOH\\
                KS-19&39&7/14/21&1/1/0&6&---&100&15&900&AS\\
                KS-19&39&7/14/21&\highlight{2/1/0}&6&---&100&15&900&MP\\
                \midrule
                \multicolumn{10}{c}{Empirical Model Values}\\
                \midrule
                \addlinespace
                \multicolumn{3}{l}{Base hit rolls depend on rate of fire}&1/1/0&&&&15\\
                &&&2/2/1&&&&120\\
                &&&3/2/1&&&&500\\
                &&&4/3/2&&&&1000\\
                &&&4/4/3&&&&1500\\
                &&&5/4/3&&&&2000\\
                \addlinespace
                \multicolumn{3}{l}{Hit roll modifiers for FCR/LRF}&\plus{1}&&R\\
                &&&\plus{1}&&L\\
                &&&\plus{2}&&W\\
                \addlinespace
                \multicolumn{3}{l}{Hit roll modifiers for proximity fuzes}&\plus{1}\\
                \addlinespace
                \multicolumn{3}{l}{Hit roll modifiers for section}&\minus{1}\\
                \addlinespace
                \multicolumn{3}{l}{Base damage depends on caliber}&&1&&\lessthanorequaltobefore{14.5}\\
                &&&&2&&20--23\\
                &&&&3&&25--30\\
                &&&&4&&35--40\\
                &&&&5&&57\\
                &&&&6&&75--100\\
                &&&&7&&\greaterthanorequaltobefore{120}\\
                \addlinespace
                \multicolumn{3}{l}{Damage modifier for rate of fire}&&\plus{1}&&&\greaterthanorequaltobefore{1800}\\
                \addlinespace
                \midrule
                \multicolumn{10}{c}{Modified or Adopted Values}\\
                \midrule
                ZPU-2&&&4/3/2&1&&&\multicolumn{3}{l}{--- adopted from AS}\\
                ZPU-4&&&5/4/3&&&&\multicolumn{3}{l}{--- for rate of fire}\\
                %Single Rh-202&&&4/3/2&&&&\multicolumn{3}{l}{--- for rate of fire}\\
                Twin Rh-202&&&5/4/3&3&&&\multicolumn{3}{l}{--- for rate of fire}\\
                ZU-23&&&5/4/3&3&&&\multicolumn{3}{l}{--- for rate of fire}\\
                Oerlikon GDF-001&&&4/3/2&&&&\multicolumn{3}{l}{--- for rate of fire}\\
                BOFI&&&5/5/3&&&&\multicolumn{3}{l}{--- for FCR-W}\\
                KS-12&&&1/1/0&&&&\multicolumn{3}{l}{--- for rate of fire}\\
                KS-19&&&1/1/0&&&&\multicolumn{3}{l}{--- for rate of fire}\\
                \midrule
                \multicolumn{10}{c}{New Values}\\
                \midrule
                M51&30&6/12/18&4/4/3&5&W&\phantom{0}75&45&854\\
                Flab Kan 38&24&6/12/18&2/1/0&5&---&\phantom{0}75&25&700\\
                90 mm M2&32&7/14/21&3/2/1&6&---&\phantom{0}90&24&820\\
                QF 3.7-inch Mk VI&45&9/18/27&3/2/1&6&--&\phantom{0}94&20&1044\\
                120 mm M1&57&11/22/33&2/2/1&7&--&120&12&945\\
                KS-30&64&12/24/36&1/1/0&7&---&130&12&970\\
                QF 5.25-inch&46&12/24/36&1/1/0&7&---&133&10&850\\
                \bottomrule
            \end{tabularx}
            \medskip
        \end{center}
    \end{minipage}

    Qualitatively, altitude and range increase with caliber, hit rolls increase with rate of fire, and damage rating increases with both caliber and rate of fire.
    This is as expected.
    Quantitatively:
    \begin{itemize}
        \item Approximately doubling the rate of fire increases the hit rolls by one, with typical values of 1/1/0 at about 15 rounds per minutes, 2/2/1 at about 150 rounds per minute, 3/2/1 at about 500 rounds per minute, 4/3/2 at about 1000, 4/4/3 at about 1500, and 5/4/2 at about 2000 and above. Intermediate rates of fire can have one of these hit rolls increased or decreased by one.
        \item FCR-W and FCR-R effectively modify the hit rolls by \plus{1} and \plus{2}.
        \item A laser rangefinder (LRF) effectively modifiers the hit rolls by \plus{1}.
        \item Proximity fuzes increase the hit rolls by \plus{1} (e.g., BOFI).
        \item As we will see below, sections have their hit rolls modified by \minus{1}.
        \item As we shall see below, sections have their hit rolls limited to 5/4/3.
        \item The base damage rating is 1 for \lessthanorequaltobefore{14.5}~mm, 2 for 20--23~mm, 3 for 25--30~mm, 4 for 35--40~mm, 5 for 57~mm, and 6 for 85--100~mm. I have extended this to 7 for \greaterthanorequaltobefore{120}~mm.
        \item The damage rating is increased by one if the rate of fire is 1800 rounds per minute or more (e.g., M55, ZPU-4, Gemag, and M167).
    \end{itemize}
    These model values are shown in the second section of the table.

    Most of the original values can be used as they are, but a few of the hit rolls and damage ratings do not fit this general pattern and need adjusting. Also, for guns for which the sources give different values, a single set needs to be adopted. These are shown in the third section of the table. Specifically:
    \begin{itemize}

        \item ZPU-2: See the note on this below.

        \item ZPU-4: See the note on this below.

        \item ZU-23: See the note on this below.

        \item Single and Twin RH-202: See the note on these below.

        \item Oerlikon GDF: The Oerlikon GDF has unexpectedly high original hit rolls of 5/4/3 for its rate of fire. The values of 4/3/2 for the ZPU-2, which has a similar rate of fire, are probably more appropriate.

        \item BOFI: The BOFI has unexpectedly low original hit rolls of 4/4/3. Values of 5/5/3, obtained by applying a \plus{2} modifier for FCR-W to the Bofors L/70, are probably more appropriate. (Note that the “BOFI” in {\AirStr} corresponds to what I call the “BOFI-R” below.)

        \item KS-19: {\MP} gives hit rolls of 2/1/0 whereas AS gives 1/1/0. While it makes sense that the hit roll is less at short range than medium range, it makes less sense that in this case the hit roll for the KS-19 (15 rounds per minute) is equal to that of the S-60 (110 rounds per minute). I prefer the {\AirStr} values.

    \end{itemize}

    Also, I note the following:

    \begin{itemize}

        \item FCR-R: The M167 and Gemag have hit rolls of 6/5/4, and the twin Rh-202 has 5/4/3. Otherwise, they are similar, except that the M167 and Gemag have FCR-R, and the Rh-202 does not (and the M167 has a moderate advantage in rate of fire). I consider the FCR-R of the M167 and Gemag to be largely responsible for increasing their hit rolls by \plus{1} compared to the Rh-202.

        \item Bofors: The hit rolls of 2/2/1 for the Bofors L/60, 3/2/1 for the 61-K (M-38), and 3/3/2 for the Bofors L/70 seem to reflect the increase in rate of fire. The L/60 and L/70 have the same maximum altitude of 16 and maximum range of 12, despite the L/60 having lower muzzle velocity. However, I am not especially motivated to attempt to change this, since the Bofors is primarily a weapon for point-defense against aircraft up to about level 10. The real BOFI and BOFI-R also have proximity fuzes, and I consider these will increase their hit rolls by \plus{1}.

        \item Proximity Fuzes: For guns with proximity fuzes (BOFI, 90 mm M2, QF 3.7-inch, and 120 mm M1), I increase the hit rolls by \plus{1}.

    \end{itemize}

    Values for new guns are shown in the fourth section of the table. They are based on the empirical model discussed here and are explained in more detail in the sections for the specific guns.

    \noteparagraph{AAA Values for Sections} I now turn to sections. Again, my aim is to empirically understand the underlying model sufficiently to identify anomalous values and to apply it to other guns.

    The upper part of the following table shows the original data for gun sections from {\AirStr} (AS), {\EOTG} (EOTG), and {\TSOH} (TSOH), and \MP' 2026 compilation (MP). Again, I have added columns with the caliber, burst rate of fire per mount, and muzzle velocity. Values in red will be subsequently modified.

    \begin{center}
        \medskip
        \smallskip
        AAA VALUES FOR SECTIONS\par
        \medskip
        \scriptsize
        \begin{tabularx}{\linewidth}{LrlllllrrL}
            \toprule
            Type
            &\multicolumn{1}{c}{\vertical{Maximum Altitude}}
            &\multicolumn{1}{c}{\vertical{Ranges}}
            &\multicolumn{1}{c}{\vertical{Hit Rolls}}
            &\multicolumn{1}{c}{\vertical{Damage}}
            &\multicolumn{1}{c}{\vertical{FCR/LRF}}
            &\multicolumn{1}{c}{\vertical{Caliber (mm)}}
            &\multicolumn{1}{c}{\vertical{Rate of fire (RPM)}}
            &\multicolumn{1}{c}{\vertical{Muzzle Velocity (m/s)}}
            &\multicolumn{1}{c}{\vertical{Source}}
            \\
            \midrule
            \multicolumn{10}{c}{Original Values}\\
            \midrule
            M16&5&2/3/4&4/3/2&2&---&12.7&2000&890&EOTG/MP\\
            Mobile ZPU-4&6&2/3/5&\highlight{3/3/2}&2&---&14.5&2400&1000&AS\\
            Mobile ZPU-4&6&2/3/5&\highlight{5/4/3}&2&---&14.5&2400&1000&MP\\
            M163&8&2/4/6&5/4/3&3&R&20&3000&1030&AS\\
            M163&8&2/4/6&\highlight{6/5/4}&3&R&20&3000&1030&MP\\
            M3 TCM-20&8&2/4/6&3/3/2&2&---&20&1450&860&EOTG\\
            Panhard&8&2/4/6&\highlight{2/2/1}&2&---&20&1450&860&AS\\
            Panhard&8&2/4/6&3/3/2&2&---&20&1450&860&MP\\
            Mobile ZU-23&9&2/5/8&\highlight{2/2/1}&2&---&23&2000&980&AS/MP\\
            Nile 23&9&2/5/8&\highlight{3/3/2}&2&---&23&2000&980&EOTG/MP\\
            ZSU-23-4&9&2/5/8&5/4/3&3&W&23&4000&980&AS/MP\\
            Blazer&10&3/6/9&5/4/3&4&W&25&1800&1040&AS\\
            Blazer&\highlight{9}&3/6/9&5/4/3&4&W&25&1800&1040&MP\\
            M53/59&11&4/7/9&\highlight{4/3/2}&3&---&30&1000&1000&AS/MP\\
            AMX-30&11&4/7/9&5/4/3&3&W&30&1200&970&AS\\
            ZSU-30-2&11&4/7/9&5/4/3&4&W&30&5000&960&AS/MP\\
            Gepard&12&4/8/10&5/4/3&4&W&35&2200&1175&AS\\
            Gepard&12&\highlight{4/8/12}&5/4/3&4&W&35&2200&1175&MP\\
            ZSU-57-2&18&5/10/15&2/1/1&5&---&57&240&1000&EOTG/MP\\
            \midrule
            \multicolumn{10}{c}{Modified or Adopted Values}\\
            \midrule
            Mobile ZPU-4&&&4/3/2&&&&\multicolumn{3}{l}{--- to match the M16 and ZPU-4}\\
            M163&&&5/4/3&&&&\multicolumn{3}{l}{--- to match the M167}\\
            Panhard&&&4/4/&&R&&\multicolumn{3}{l}{--- to match the M3 TCM-20}\\
            Mobile ZU-23&&&4/3/2&3&&&\multicolumn{3}{l}{--- to match the ZU-23}\\
            Nile 23&&&4/3/2&3&&&\multicolumn{3}{l}{--- to match the ZU-23}\\
            Blazer&10&&&&&&\multicolumn{3}{l}{--- adopt {\AirStr} values}\\
            M53/59&&&3/2/1&&&&\multicolumn{3}{l}{--- to match the AMX-30}\\
            Gepard&&4/8/10&&&&&\multicolumn{3}{l}{--- adopt {\AirStr} values}\\
            \midrule
            \multicolumn{10}{c}{New Values}\\
            \midrule
            AMX-13 DCA&11&4/7/9&5/4/3&3&W&30&1200&970&\\
            M15&12&2/4/8&3/2/1&4&---&37&120&792\\
            M19&15&4/8/12&2/2/1&4&---&40&270&865&\\
            M34&15&2/4/82&1/1/1&4&---&40&135&865&\\
            M42&15&4/8/12&2/2/1&4&---&40&270&865&\\
            \bottomrule
        \end{tabularx}
        \medskip
    \end{center}

    Comparing similar weapons, batteries and sections have the same maximum altitudes, ranges, and damage ratings.
    However, sections (with two mounts) typically have a \minus{1} modification to their hit rolls compared to batteries (with four to six mounts).
    In a very few cases, the difference is 0, and in a very few others, it is \minus{2}.
    A difference of 1 in the hit rolls is commensurate with the approximate doubling of the rate of fire from a section to a battery.

    To determine the appropriate AAA for sections, the following recipe seems to work in most cases:
    \begin{itemize}
        \item Take the value for an equivalent battery apply an adjustment of \minus{1} to the hit rolls to account for the reduced rate of fire in the section. (Only do this if the rate of fire is reduced; see the comment on the ZSU-23-4 and ZSU-57-2 below.)
        \item Limit the hit rolls to 5/4/3. The preceding guideline would predict hit rolls of 6/5/4 for the Blazer and 7/6/5 for the ZSU-23-4 and ZSU-30-2, so this limit is to conform to the original values. I suggest that in the underlying model, beyond a certain rate of fire, other factors take over to limit the probability of a hit.
    \end{itemize}

    Again, most of the original data can be used as they are, but a few of the values are worth commenting upon or need adjusting.
    \begin{itemize}

        \item Mobile ZPU-4: See the note on this below.

        \item M163: See the note on this below.

        \item Panhard M3 VDA: See the note on this below.

        \item Mobile ZU-23 and Nile 23: The original hit rolls for the mobile ZU-23 and Nile 23 section are 2/2/1 and 3/3/2, despite them being quite similar. The model values of 4/3/2 are probably more appropriate and are in better agreement with the values for the M16. Also, as with the ZU-23 battery, a damage rating of 3 is probably more appropriate given their high rate of fire.

        \item Blazer: {\MP} gives this a maximum altitude of 9 rather than the 10 in {\AirStr}. The Blazer uses a 25~mm gun, and the {\AirStr} value is between that of 23~mm guns (9) and 30~mm guns (11), so I will adopt it. Following the model, the hit rolls would be 6/5/4, but they are limited to 5/4/3.

        \item ZSU-23-4: The ZSU-23-4 has four barrels whereas the ZU-23-2 has two, so the normal \minus{1} modifier for sections does not apply. Following the model and with the \plus{2} modifier for FCR-W, the hit rolls would be 7/6/5, but they are limited to 5/4/3.

        \item ZSU-30-2: Following the model, the hit rolls would be 7/6/5, but they are limited to 5/4/3.

        \item M53/59: See the note on this below.

        \item Gepard: {\MP} increases the long range to 12, from 10 in {\AirStr}. This might be correct; perhaps later ammunition give longer range than the GDF-001 and I note that he gives a similar increase to the GDF-005. However, I can find no evidence for such an improvement. I adopt the {\AirStr} value.

        \item ZSU-57-2: The ZSU-57-2 section has the same hit rolls as the S-60 battery. This is probably appropriate, as the ZSU-57-2 is a twin mount and this largely compensates for the reduced rate of fire of the section.

    \end{itemize}

    I adopt these modified values.

    Also, I note:

    \begin{itemize}

        \item Panhard M3 VDA: The Panhard M3 VDA has an FCR-R capability, which is discussed below but not considered here.

        \item AMX-30 DCA: The real AMX-30 DCA does not have an FCR. The AMX-30 DCA in the game seems to correspond to the similar AMX-30SA, which does.

    \end{itemize}

    The AAA values for new guns are discussed when they are presented.

}

All AAA platoons and batteries are also available as sections (but not squads), with damage resiliences of 2, hit rolls reduced by 1, and VPs corresponding to a platoon or battery that has taken 1D damage.

% \subsection{FCRs}
%
% Some AAA units have integrated FCRs that give them an all-weather and night capability.
% Most medium and heavy can also use add-on FCRs, described below.
% Units without an FCR (or with a range-only FCR) are limited to visually sighted targets.

\section{Light AAA Units}

\begin{withcounters}

    \subsectionwithcounter{M2 platoon}{M2 Platoon}

    The Browning M2 .50~cal (12.7~mm) heavy machine gun is normally found on a low tripod for ground combat or mounted on armored fighting vehicles, but here is used on a tall dual-purpose mount suitable for air defense. A platoon has six guns.
    The M2 gun was introduced in 1933 and has been used by United States forces and their allies in numerous conflicts since then.%
    \note{M2 Platoon}{

        \aaamountvalues{12.7}{500}{890}

        \input{aaavalues/M2 platoon.tex}

        I adopt the AAA values, defense strength, and sighting range from {\AirStr} and {\MP}.
        The adopted AAA values are in agreement with the empirical model.

        \begin{notesources}
            \item Wikipedia on the \href{https://en.wikipedia.org/wiki/M2_Browning}{M2}.
            \item WeaponSystems on the \href{https://weaponsystems.net/system/411-Browning+M2HB}{M2}.
        \end{notesources}
    }

    \subsectionwithcounter{DShK platoon}{DShK Platoon}

    The DShK M1938 12.7~mm heavy machine gun is most commonly found on a low, wheeled mount for ground combat or mounted on armored fighting vehicles, but here is used on a tall tripod suitable for air defense. It is nicknamed the Dushka (darling) or Dochka (daughter). A platoon has six guns.
    The DShK entered service with the Soviet Army in 1938 and saw action in WW2 and numerous conflicts since then, including the Korean and Vietnam Wars, in which the Soviet Union participated or provided arms. In particular, it was the main air-defense gun for Chinese and North Korean field units in the early part of the Korean War.%
    \note{DShK Platoon}{

        \aaamountvalues{12.7}{500}{850}

        \input{aaavalues/DShK platoon.tex}

        I adopt the AAA values, defense strength, and sighting range from {\AirStr} and {\MP}.
        The adopted AAA values are in agreement with the empirical model.

        Service in the Korean War is from Singh, ch.~1.

        % LTeX: dictionary+=Defence
        \begin{notesources}
            \item Wikipedia on the \href{https://en.wikipedia.org/wiki/DShK}{DShK}.
            \item WeaponSystems on the \href{https://weaponsystems.net/system/744-DShK}{DShK}.
            \item Singh, M. (Pen \& Sword, 2020): “Anti-Aircraft Artillery in Combat 1950--1972: Air Defence in the Jet Age”.
        \end{notesources}
        % LTeX: dictionary-=Defence

    }

    \subsectionwithcounter{M16 section}{M16 Section}

    The M16 Multiple Gun Motor Carriage has an M45 turret with four M2 .50 cal (12.7~mm) heavy machine guns mounted on a modified M3 half-track. A section has two vehicles. The M16 is a development on the M13 and M14 Multiple Gun Motor Carriages, which have two M2 heavy machine guns in a similar turret. As a single M2 has an insufficient rate of fire to be effective even against WW2 aircraft, first two and then four were used on a single mount.
    The M16 was used by United States forces during WW2. It was also used by US, Commonwealth, and South Korean forces in the Korean War, primarily in an anti-infantry role. It was largely withdrawn form US services in late 1951 and officially retired in 1958. Israeli forces employed it during the 1956 Arab-Israeli War. Other post-WW2 users include Belgium, France, West Germany, Japan, Netherlands, Philippines, Portugal, and Thailand%
    \note{M16 Section}{

        \aaamountvalues{12.7}{2000}{890}

        \input{aaavalues/M16 section.tex}

        I adopt the AAA values, defense strength, and sighting range from {\EOTG} and {\MP}.
        The adopted AAA values are in agreement with the empirical model.

        \begin{notesources}
            \item Wikipedia on the \href{https://en.wikipedia.org/wiki/M16_Multiple_Gun_Motor_Carriage}{M16}, \href{https://en.wikipedia.org/wiki/M45_quad_mount}{M45}, and \href{https://en.wikipedia.org/wiki/M2_Browning}{M2}.
            \item WeaponSystems on the \href{https://weaponsystems.net/system/739-M16\%20MGMC}{M16} and \href{https://weaponsystems.net/system/411-Browning+M2HB}{M2}.
        \end{notesources}
    }

    \subsectionwithcounter{M17 section}{M17 Section}

    The M17 Multiple Gun Motor Carriage is the equivalent of the M16 Multiple Gun Motor Carriage, but based on the slightly different M5 half-track.
    A section has two vehicles.
    The M17 was supplied to the Soviet Union and entered service in 1944. It was presumably replaced by the BTR-40A.%
    \note{M17 Section}{

        \aaamountvalues{12.7}{2000}{890}

        \input{aaavalues/M17 section.tex}

        I give this same AAA values, defense strength, and sighting range as the M16.
        The adopted AAA values are in agreement with the empirical model.

        \begin{notesources}
            \item Wikipedia on the \href{https://en.wikipedia.org/wiki/M16_Multiple_Gun_Motor_Carriage}{M17}, \href{https://en.wikipedia.org/wiki/M45_quad_mount}{M45}, and \href{https://en.wikipedia.org/wiki/M2_Browning}{M2}.
            \item Tank Encyclopedia on the \href{https://tanks-encyclopedia.com/ww2/usa/multiple-gun-motor-carriage-m16-and-m17/}{M17}.
        \end{notesources}
    }

    \subsectionwithcounter{M55 platoon}{M55 Platoon}

    The M55 has an M45 turret with four M2 .50 cal (12.7~mm) heavy machine guns (the same turret as used in the M16) mounted on a towed carriage. A platoon has four gun mounts.
    The M55 was used by United States forces during the Korean War and by Pakistani forces during the 1965 India-Pakistan War.%
    \note{M55 Battery}{

        \aaamountvalues{12.7}{2000}{890}

        \input{aaavalues/M55 platoon.tex}

        I adopt the AAA values, defense strength, and sighting range from {\EOTG} and {\MP}.
        The adopted AAA values are in agreement with the empirical model.

        \begin{notesources}
            \item Wikipedia on the \href{https://en.wikipedia.org/wiki/M45_quad_mount}{M45} and \href{https://en.wikipedia.org/wiki/M2_Browning}{M2}.
            \item WeaponSystems on the \href{https://weaponsystems.net/system/411-Browning+M2HB}{M2}.
        \end{notesources}
    }

    \subsectionwithcounter{Vz.53 platoon}{Vz.53 Platoon}

    The Czechoslovak Vz.53 has four Vz.38/46 12.7 mm machine guns (copies of the Soviet DShK) on a light, wheeled mount.
    It is sometimes referred to as the M53.
    A platoon has six guns.
    It entered service with the Czechoslovak Army in 1953, and was exported to Cuba, Egypt (from 1955), Iraq, Syria, and Vietnam. Cuba subsequently exported some to Grenada. It saw action with the Egyptian Army in the 1956 War, with Cuba in the Bay of Pigs Invasion in 1961, in the Vietnam War, in the Nicaraguan Revolution, with the Grenadian Army during the invasion by US force in 1983%
    \note{Vz.53 Platoon}{

        \aaamountvalues{12.7}{2000}{850}

        \input{aaavalues/Vz.53 platoon.tex}

        I give this same AAA values, defense strength, and sighting range as the M55.
        The adopted AAA values are in agreement with the empirical model.

        {\MP} gives data for the M53, but it appears to be a truck-mounted section (the “type” is given as “truck” and the hit rolls are the same as an M16 section), but one inconsistent value is its defense strength of 3, which corresponds to other towed batteries, rather than 2, which corresponds to other truck-mounted sections.

        \begin{notesources}
            \item Wikipedia on the \href{https://en.wikipedia.org/wiki/Vz.53_anti-aircraft_gun}{Vs.53} and \href{https://en.wikipedia.org/wiki/DShK}{DShK}.
            \item WeapoNEWS on the \href{https://weaponews.com/weapons/65350809-post-war-anti-aircraft-guns-of-czechoslovakia.html}{Vz.53}.
            \item WeaponSystems on the \href{https://weaponsystems.net/system/744-DShK}{DShK}.
        \end{notesources}
    }

    \subsectionwithcounter{BTR-152A Vz.53 section}{BTR-152A Vz.53 Section}

    The BTR-152A Vz.53 has a Czechoslovak Vz.53 gun mount with four 12.7~mm Vz.38/46 machine guns in the open, six-wheeled Soviet BTR-152 APC.
    The gun mount is also referred to as the M53.
    A section has two vehicles.
    It served only in the Egyptian Army, and saw action in the 1956 War.%
    \note{BTR-152A Vz.53 Section}{

        \aaamountvalues{12.7}{2000}{850}

        \input{aaavalues/BTR-152A Vz.53 section.tex}

        I give this same AAA values, defense strength, and sighting range as the M16.
        The adopted AAA values are in agreement with the empirical model.

        The in-service date assumes they were supplied or converted shortly after the 1955 arms deal.
        Some {\EOTG} scenarios for 1956 War give the Egyptian Army M16s, but I do not think these served with the Egyptian Army; I think these might well be BTR.152A Vz.53s.

        The BTR-152 is based on a truck, and so I give it “G/P” mobility rather than “U”.

        \begin{notesources}
            \item Wikipedia on the \href{https://en.wikipedia.org/wiki/BTR-152\#Soviet_Union}{BTR-152A}, \href{https://en.wikipedia.org/wiki/Vz.53_anti-aircraft_gun}{Vs.53}, and \href{https://en.wikipedia.org/wiki/DShK}{DShK}.
            \item WeapoNEWS on the \href{https://weaponews.com/weapons/65350809-post-war-anti-aircraft-guns-of-czechoslovakia.html}{Vz.53}.
            \item Army Guide on the \href{https://www.army-guide.com/eng/product1120.html}{BTR-152}.
            \item Miranda, M., 2015: “\href{https://thesovietarmourblog.blogspot.com/2015/12/}{Gun Truck}”
        \end{notesources}
    }

    \subsectionwithcounter{mobile M45 section}{Mobile M45 Section}

    The mounted M45 has the M45 turret with four M2 .50 cal (12.7~mm) heavy machine guns (the same turret as used in the M16) removed from an M55 trailer and mounted on a truck.
    During the Vietnam War, in 1967, US forces improvised these \emph{gun trucks} for convoy security%
    \note{Mobile M45 Battery}{

        \aaamountvalues{12.7}{2000}{890}

        \input{aaavalues/mobile M45 section.tex}

        I give this same AAA values as the M16.
        The adopted AAA values are in agreement with the empirical model.
        I give this the same defense strength and sighting range as the mobile ZPU-4 in {\AirStr}.

        \begin{notesources}
            \item Wikipedia on the \href{https://en.wikipedia.org/wiki/M45_quad_mount}{M45}, \href{https://en.wikipedia.org/wiki/M2_Browning}{M2}, and \href{https://en.wikipedia.org/wiki/Gun_truck
                }{gun trucks}.
            \item WeaponSystems on the \href{https://weaponsystems.net/system/411-Browning+M2HB}{M2}.
        \end{notesources}
    }

    \subsectionwiththreecounters{ZPU-1 battery}{ZPU-2 battery}{ZPU-4 battery}{ZPU-1/2/4 Battery}

    The ZPU-1/2/4 have one, two, or four 14.5~mm heavy machine guns in open mounts on towed carriages. A battery has six gun mounts.
    The ZPU-1/2/4 was introduced in 1949 and saw extensive service with Soviet or Soviet-supported forces in many subsequent conflicts, including with communist forces in the Korean War and Vietnam War and with Arab forces in the various Arab-Israeli Wars from 1956 onwards.%
    \note{ZPU-1/2/4 Battery}{

        \aaamountvalues{14.5}{600/1200/2400}{1000}

        ZPU-1 \input{aaavalues/ZPU-1 battery.tex}
        ZPU-2 \input{aaavalues/ZPU-2 battery.tex}
        ZPU-4 \aaavalues{6}{2/3/5}{5/4/3}{2}{}


        For the ZPU-1, I adopt the AAA values, defense strength, and sighting range from {\EOTG}, {\TSOH}, and {\MP}.
        The adopted AAA values are in agreement with the empirical model.

        For the ZPU-2, {\AirStr} gives hit rolls of 4/3/2 whereas {\MP} gives it 4/3/1 and {\AirStr} gives a damage rating of 1 (equal to the ZPU-1) whereas {\MP} gives it a rating of 3 (more than the ZPU-4). The ZPU-2 should logically have both hit rolls and damage ratings intermediate to those of the ZPU-1 and ZPU-4, but this is impossible as these only differ by one. The solution in {\AirStr} seems to be a good compromise: give it the hit rolls of the ZPU-4 but the damage rating of the ZPU-1. I adopt the AAA values, defense strength, and sighting range from {\AirStr}.
        The adopted AAA values are in agreement with the empirical model.

        For the ZPU-4, {\AirStr} and {\MP} give hit rolls of 5/4/3 whereas {\TSOH} gives 4/3/2. All give damage ratings of 2. I see no reason to give this different hit rolls to those of the M55, given the similarities between the two in caliber, rate of fire, and muzzle velocity. Therefore, I give it hit rolls of 5/4/3, by analogy with the M55 and in agreement with the empirical model. I therefore adopt the AAA values, defense strength, and sighting range from {\AirStr} and {\MP}.

        \begin{notesources}
            \item Wikipedia on the \href{https://en.wikipedia.org/wiki/ZPU}{ZPU}.
            \item WeaponSystems on the \href{https://weaponsystems.net/system/1674-ZPU-1}{ZPU-1}, \href{https://weaponsystems.net/system/935-ZPU-2}{ZPU-2}, and \href{https://weaponsystems.net/system/934-ZPU-4}{ZPU-4}, and \href{https://weaponsystems.net/system/356-Vladimirov\%20KPV}{KPV}.
        \end{notesources}
    }

    \vspacelines{3}

    \subsectionwiththreecounters{mobile ZPU-1 section}{mobile ZPU-2 section}{mobile ZPU-4 section}{Mobile ZPU-1/2/4 Section}

    The mobile ZPU-1/2/4 is a light truck with a ZPU-1/2/4 gun mount.
    A section has two vehicles.
    Many users of the ZPU-1/2/4 adapt it in this way for increased mobility, and
    like the towed versions, these mounted versions have seen combat in many conflicts.
    Formal conversions tend to be used for air-defense, but informal conversions or \emph{technicals} are mainly used for ground combat, serving as light cavalry or infantry support.%
    \note{Mobile ZPU-1/2/4 Section}{

        \aaamountvalues{14.5}{600/1200/2400}{1000}

        Mobile ZPU-1 \aaavalues{6}{2/3/5}{2/1/0}{1}{}

        Mobile ZPU-2 \input{aaavalues/mobile ZPU-2 section.tex}
        Mobile ZPU-4 \input{aaavalues/mobile ZPU-4 section.tex}

        For the mobile  ZPU-1 and ZPU-2 sections, I adopt the AAA values of the ZPU-1 and ZPU-2 batteries with a \minus{1} modifier to the hit rolls for being sections. I adopt same the defense strength and sighting range as the mobile ZPU-4 section in {\AirStr} and {\MP}.
        The adopted AAA values are in agreement with the empirical model.

        For the mobile ZPU-4, {\AirStr} gives hit rolls of 3/3/2 and {\MP} gives hit rolls of {5/4/3}. The values from {\AirStr} seem to be too low and those of {\MP} seem to be too high. I see no reason to give this different hit rolls to those of the M16, given the similarities between the two in caliber, rate of fire, and muzzle velocity. Therefore, I give it hit rolls of 4/3/2, by analogy with the M16, in agreement with those of the ZPU-4 battery with a \minus{1} modifier, and in agreement with the empirical model.
        I adopt the defense strength and sighting range from {\AirStr} and {\MP}.

        %{\AirStr} gives this 5 VPs, which is the same as the towed ZPU-4 battery despite the mobile section having worse hit rolls and defense strength. I think 4 VPs (i.e., one less than the battery) is more appropriate.

        \begin{notesources}
            \item Wikipedia on the \href{https://en.wikipedia.org/wiki/ZPU}{ZPU} and \href{https://en.wikipedia.org/wiki/Technical_(vehicle)}{Technicals}.
            \item WeaponSystems on the \href{https://weaponsystems.net/system/1674-ZPU-1}{ZPU-1}, \href{https://weaponsystems.net/system/935-ZPU-2}{ZPU-2}, and \href{https://weaponsystems.net/system/934-ZPU-4}{ZPU-4}, and \href{https://weaponsystems.net/system/356-Vladimirov\%20KPV}{KPV}.
        \end{notesources}
    }

    \vspacelines{3}

    \subsectionwithcounter{BTR-40A section}{BTR-40A Section}

    The BTR-40A has a twin 14.5~mm ZPU-2 turret in the open, four-wheeled BTR-40 APC.
    A section has two vehicles.
    The BTR-40A entered service with the Soviet Army in 1950, presumably replacing the M17.
    It was also used by East Germany and Cuba.%
    \note{BTR-40A Section}{

        \aaamountvalues{14.5}{1200}{1000}

        \input{aaavalues/BTR-40A section.tex}

        I have adopted the AAA values of the mobile ZPU-2 section and the defense strength and sighting range of the M16.
        The adopted AAA values are in agreement with the empirical model.

        The BTR-40 is based on a truck, and so I give it “G/P” mobility rather than “U”.

        \begin{notesources}
            \item Wikipedia on the \href{https://en.wikipedia.org/wiki/BTR-40\#Soviet_Union_2}{BTR-40A} and \href{https://en.wikipedia.org/wiki/ZPU}{ZPU-2}.
        \end{notesources}
    }

    \subsectionwithtwocounters{BTR-152A ZPU-2 section}{BTR-152A ZPU-4 section}{BTR-152A ZPU-2/4 Section}

    The BTR-152A has a twin 14.5~mm ZPU-2 or quad 14.5~mm ZPU-4 turret in the open, six-wheeled BTR-152 APC.
    A section has two vehicles.
    The BTR-152A entered service with the Soviet Army in 1951, presumably replacing the BTR-40A.
    The BTR-152D/E are similar, but based on the improved BTR-152V/V1 variants.
    Other users include Ethiopia in the Ogaden War and Egypt.%
    \note{BTR-152A Section}{

        \aaamountvalues{14.5}{1200/2400}{1000}

        BTR-152A ZPU-2 \aaavalues{6}{2/3/5}{3/2/1}{1}{}

        BTR-152A ZPU-4 \input{aaavalues/BTR-152A ZPU-4 section.tex}

        I have adopted the AAA values of the mobile ZPU-2/4 section and the defense strength and sighting range of the M16.
        The adopted AAA values are in agreement with the empirical model.

        The BTR-152 is based on a truck, and so I give it “G/P” mobility rather than “U”.

        Army Guide states that the A/E mounted a ZPU-2 and the D a ZPU-4. Others (Wikipedia) that the A/D/E referred to the chassis (with the D using the BTR-152V and the E the BTR-152V1 chassis) rather than the gun mount. I am not sure how to resolve this, but I don't think it is that important. Instead, I will use “BTR-152A ZPU-2” and “BTR-152A ZPU-4,” which extends obviously to other guns mounted on the BTR-152A and similar.

        \begin{notesources}
            \item Wikipedia on the \href{https://en.wikipedia.org/wiki/BTR-152\#Soviet_Union}{BTR-152A} and \href{https://en.wikipedia.org/wiki/ZPU}{ZPU-2}.
            \item Army Guide on the \href{https://www.army-guide.com/eng/product1120.html}{BTR-152}.
            \item Miranda, M., 2015: “\href{https://thesovietarmourblog.blogspot.com/2015/12/}{Gun Truck}”
        \end{notesources}
    }

    \subsectionwithcounter{M167 platoon}{M167 Platoon}

    The M167 Vulcan Air Defense System (VADS) has a 20~mm six-barreled Vulcan gun in an open mount on a towed carriage. A platoon has four gun mounts. The fire-control system has optical sights and a range-only radar.
    The M167 was introduced into service with the US Army in 1967, principally for airfield defense, and subsequently served with other US allies.%
    \note{M167 Platoon}{

        \aaamountvalues{20}{3000}{1030}

        \aaavalues{8}{2/4/6}{6/5/4}{3}{R/HF}


        I have adopted the AAA values, defense strength, and sighting range from {\AirStr} and {\MP}.
        The adopted AAA values are in agreement with the empirical model.

        The 1976 Operator's Manual describes the use of the AN/TVS-2B or AN/TVS-5 night sight. These are first- and second-generation image intensifiers. They are mounted next to the tracking sight and appear to be useful only for ground targets. I do not count it as giving “night IR sights” capability.

        MobileRadar states the AN/VPS-2 operates at 9.2--9.5 GHz.

        \begin{notesources}
            \item Wikipedia on the \href{https://en.wikipedia.org/wiki/M167_VADS}{M167} and \href{https://en.wikipedia.org/wiki/M61_Vulcan}{M61 Vulcan}.
            \item WeaponSystems on the \href{https://weaponsystems.net/system/150-M167+Vulcan}{M167} and \href{https://weaponsystems.net/system/1510-20mm\%20M61\%20Vulcan}{M61}.
            \item Department of the Army, \href{https://archive.org/details/TM9-1005-286-10/mode/2up}{TM 9-1005-286-10}, 1976: “Operator's Manual (Crew) for Gun, Air Defense Artillery, Towed”.
            \item \href{https://www.mobileradar.org/radar_descptn_3.html}{MobileRadar} on the AN/VPS-2.
        \end{notesources}
    }

    \subsectionwithcounter{M163 section}{M163 Section}

    The M163 Vulcan Air Defense System (VADS) has a 20~mm six-barreled Vulcan gun in a lightly armored but open-topped turret on an adapted M113 APC. A section has two vehicles. The fire-control system has optical sights and a range-only radar.
    The M163 VADS was introduced into service with the US Army in 1968 and has since been adopted by other US allies. The M163 was used by US forces in the Vietnam War, the 1989 Invasion of Panama, and the 1991 Gulf War. It was also used by the Israel Army in the 1982 Lebanon War and in 2002 during the Second Intifada.
    From 1988 in the US Army, each M163 was typically issued with a FIM-92 Stinger launcher with two rounds, so each M163 section effectively carried a pair of FIM-92 squads.%
    \note{M163 Section}{

        \aaamountvalues{20}{3000}{1030}

        \input{aaavalues/M163 section.tex}

        I have adopted the AAA values, defense strength (4), and sighting range (12) from {\AirStr}, but open rather than hard target class.
        The adopted AAA values are in agreement with the empirical model.
        {\MP} gives the section the same hit rolls as the M167 battery, despite the fewer mounts, which does not seem to be correct.
        The turret is open, so I give this a target class open rather than hard.

        The 1969 report and 1976 manual mention the AN/VPS-2 range-only radar. MobileRadar states it operates at 9.2--9.5 GHz.

        The 1976 Operator's Manual describes the use of the AN/TVS-2B night sight. This is a first-generation image intensifier. It is mounted next to the tracking sight and appears to be useful only for ground targets.  I do not count it as giving “night IR sights” capability.

        The PIVADS version introduced in 1984 had an improved sight, but I assume this makes no difference in the AAA values.

        \begin{notesources}
            \item Wikipedia on the \href{https://en.wikipedia.org/wiki/M163_VADS}{M163}, \href{https://en.wikipedia.org/wiki/M167_VADS}{M167}, and \href{https://en.wikipedia.org/wiki/M61_Vulcan}{M61 Vulcan}.
            \item WeaponSystems on the \href{https://weaponsystems.net/system/67-M163+Vulcan}{M163}, \href{https://weaponsystems.net/system/150-M167+Vulcan}{M167} and \href{https://weaponsystems.net/system/1510-20mm\%20M61\%20Vulcan}{M61}.
            \item Department of the Army, \href{https://archive.org/details/TM9235030010M163VADS/mode/2up}{TM 9-2350-300-10}, 1976: “Operator's Manual For Gun, Air Defense Artillery, Self-Propelled 20-mm, M163”.
            \item Department of the Army, \href{https://apps.dtic.mil/sti/tr/pdf/AD0857327.pdf}{ACGC-64F}, 1969: “Final Report: XM163 Vulcan Air Defense System”.
            \item \href{https://www.mobileradar.org/radar_descptn_3.html}{MobileRadar} on the AN/VPS-2.
        \end{notesources}
    }

    \subsectionwithcounter{TCM-20 platoon}{TCM-20 Platoon}

    The TCM-20 is an Israeli adaptation of the M55, with two Hispano-Suiza 20mm guns from obsolete aircraft replacing the four .50 cal machine guns. As with the M55, the fire-control system is entirely optical. A platoon has four gun mounts.
    The TCM-20 was employed by Israeli forces from 1970 and subsequently saw combat in the 1973 Arab-Israeli War and the 1982 Lebanon War.%
    \note{TCM-20 Platoon}{

        \aaamountvalues{20}{1450}{860}

        \input{aaavalues/TCM-20 platoon.tex}

        I have adopted the AAA values, defense strength, and sighting range from {\EOTG}.
        The adopted AAA values are in agreement with the empirical model.

        \begin{notesources}
            \item Wikipedia on the \href{https://en.wikipedia.org/wiki/M45_quad_mount\#TCM-20}{TCM-20} and \href{https://en.wikipedia.org/wiki/Hispano-Suiza_HS.404}{HS.404}.
            \item WeaponSystems on the \href{https://weaponsystems.net/system/750-TCM-20}{TCM-20}.
        \end{notesources}
    }

    \subsectionwithcounter{M3 TCM-20 section}{M3 TCM-20 Section}

    The M3 TCM-20 has a TCM-20 turret Hispano-Suiza 20mm guns mounted in an adapted M3 half-track, like the earlier M16. A section has two vehicles.
    The M3 TCM-20 was used by Israeli forces in the 1973 Arab-Israeli War and the 1982 Lebanon War.%
    \note{M3 TCM-20 Section}{

        \aaamountvalues{20}{1450}{860}

        \input{aaavalues/M3 TCM-20 section.tex}

        I have adopted the AAA values, defense strength, and sighting range from {\EOTG}.
        The adopted AAA values are in agreement with the empirical model.

        \begin{notesources}
            \item Wikipedia on the \href{https://en.wikipedia.org/wiki/M3_half-track\#Post-war_Israeli_variants}{M3 TCM-20}, \href{https://en.wikipedia.org/wiki/M45_quad_mount\#TCM-20}{TCM-20}, \href{https://en.wikipedia.org/wiki/Hispano-Suiza_HS.404}{HS.404}.
            \item WeaponSystems on the \href{https://weaponsystems.net/system/786-M3\%20TCM-20}{M3 TCM-20} and \href{https://weaponsystems.net/system/750-TCM-20}{TCM-20}.
        \end{notesources}
    }

    \subsectionwithcounter{Rh-202 battery}{Rh-202 Battery}

    The Rheinmetall Rh-202 has a twin 20mm gun in an open mount on a towed carriage. The fire-control system is entirely optical.
    A battery has six gun mounts.
    It entered service with West German forces in 1972 and was used for airfield defense. Subsequently, it was exported and notably used by the Argentine Air Force in the 1982 South Atlantic War in the defense of BAM Malvinas (Port Stanley Airport) and BAM Condor (Goose Green Airfield).%
    \note{Rh-202 Battery}{

        \aaamountvalues{20}{2060}{1044}

        Rh-202 \aaavalues{8}{2/4/6}{5/4/3}{3}{}


        % For the single Rh-202, {\AirStr} gives unexpectedly low original hit rolls of 3/2/1 given its rate of fire.
        % The ZPU-2 and TCM-20 have a similar rate of fire and values of 4/3/2.
        % Therefore, by analogy with the ZPU-2 and TCM-20 and in agreement with the empirical model, I adopt AAA values, defense strength, and sighting range from {\AirStr}, but modify the hit rolls to 4/3/2.

        I adopt the AAA values, defense strength, and sighting range from {\AirStr}, except I increase the damage rating from 2 to 3.
        {\AirStr} gives an unexpectedly low original damage rating of 2 for its rate of fire. A value of 3, increased by one from the baseline for the caliber, is probably more appropriate.
        The adopted AAA values are in agreement with the empirical model.

        {\MP} gives values for the Rh-202, but it is not clear whether they are for the single or twin mount. He gives hit rolls of 3/2/1 (equal to the {\AirStr} values for the single mount) but a damage rating of 3 (equal to the value I have adopted for the twin mount). I do not adopt this combination of values either for the single or twin Rh-202.

        {\AirStr} has single and twin Rh-202 batteries.
        Single mounts seem to be used for AFVs (e.g., Marder, Luchs, and Wiesel) and naval applications, but all AAA mounts seem to be twin. I therefore drop the single mounts.

        \begin{notesources}
            \item Wikipedia on \href{https://en.wikipedia.org/wiki/Rheinmetall_Mk_20_Rh-202}{Rh-202}.
            \item WeaponSystems on the \href{https://weaponsystems.net/system/937-Rh-202}{Rh-202 twin mount} and \href{https://weaponsystems.net/system/415-20mm+Rheinmetall+Rh-202}{Rh-202 gun}.
        \end{notesources}
    }

    \subsectionwithcounter{Panhard M3 VDA section}{Panhard M3 VDA Section}

    The Panhard M3 VDA vehicles have an armored turret for twin HS.820 20~mm guns mounted on the Panhard M3 APC.
    A section has two vehicles.
    The guns have optical sights, but typically one vehicle in each section or platoon is equipped with a search-and-ranging radar and can share this information automatically with the others.
    The Panhard M3 VDA served with the armed forces of Côte d'Ivoire, Niger, and the UAE.%
    \note{Panhard M3 VDA Section}{

        \aaamountvalues{20}{1450}{860}

        \aaavalues{8}{2/4/6}{4/4/3}{2}{R/UF}


        I have adopted the AAA values, defense strength, and sighting range from {\AirStr}, except that I have adopted hit rolls of 4/4/3 and FCR-R.

        The Panhard M3 VDA has the same twin guns as the M3 TCM-20.
        {\AirStr} has unexpectedly low original hit rolls of 2/2/1 compared to those of 3/3/2 for the M3 TCM-20 section.
        {\MP} gives it the same hit rolls, but omits the FCR-R.
        I adopt the hit rolls of the M3 TCM-20 (3/3/2) with a \plus{1} modifier for FCR-R.
        The adopted AAA values are in agreement with the empirical model.

        The Panhard M3 VDA in {\AirStr} seems to correspond to Wikipedia M3 VDA or WeaponSystems M3 VDAA.

        However, the Panhard M3 VDA radar-equipped vehicle seems to be able to search and share speed and range information with the other vehicle in the section, which seems equivalent to FCR-R.

        The radar is the RA 20. Janes states that it operates in the E band. There is no mention of IFF, but it does seem to be a PD radar with a MTI.

        % LTeX: dictionary+=Defence
        \begin{notesources}
            \item Wikipedia on the \href{https://en.wikipedia.org/wiki/Panhard_M3\#Variants}{Panhard M3 VDA} and \href{https://en.wikipedia.org/wiki/Hispano-Suiza_HS.820}{HS.820}.
            \item WeaponSystems on the \href{https://weaponsystems.net/system/532-VDAA}{Panhard M3 VDAA}.
            \item Jane’s Land-Based Air Defence 1992-93 (\href{https://forum.warthunder.com/t/british-weapon-systems-technical-data-and-discussion/3879/2225}{extract})
        \end{notesources}
        % LTeX: dictionary-=Defence

    }

    \subsectionwithcounter{ZU-23 battery}{ZU-23 Battery}

    The ZU-23 has twin 23~mm cannons in an open mount on a towed carriage. It was designed as a more powerful replacement for the ZPU-1/2/4.
    A battery has six gun mounts.
    The fire-control system is entirely optical.
    The ZU-23 was introduced into service by the Soviet Army in 1960. It later served with the armed forces of many Soviet allies, including with North Vietnam in the Vietnam War, with Egypt and Syria in the 1967 and 1973 Wars, with Soviet forces in the Soviet-Afghan War, with Iraq in the Iraq-Iran War, with Syria in the 1982 Lebanon War, and with both sides in the Afghan Civil War.%
    \note{ZU-23 Battery}{

        \aaamountvalues{23}{2000}{980}

        \aaavalues{9}{2/5/8}{5/4/3}{3}{}


        I adopt the AAA values, defense strength, and sighting range are from {\AirStr}, {\TSOH}, and {\MP}, but the hit rolls are modified as from 4/3/2 to 5/3/2 and the damage rating from 2 to 3.

        {\AirStr}, {\TSOH}, and {\MP} all give the ZU-23 hit rolls of 4/3/2, but give 5/4/3 to other AAA batteries with similar rates of fire (M55, twin Rh-202, M167, and Gemag) after accounting for modifiers for FCR.
        Furthermore, all of these other batteries have damage ratings increased by one above their baseline values.
        There does not seem to be a reason to penalize the ZU-23.
        It is undeniably basic, and one might think that would be reason to reduce its effectiveness. However, the TCM-20 is basic too and has both a lower rate of fire than the ZU-23 (1450 compared to 2000) and a smaller caliber (20~mm compared to 23~mm), yet its hit rolls are higher than the ZU-23 in {\AirStr}, {\TSOH}, and {\MP} and its damage rating is the same.

        Therefore, by analogy with these batteries and in agreement with the empirical model, I give it hit rolls of 5/4/3 and a damage rating of 3.

        %TSOH gives it 6 VPs and {\AirStr} gives it only 5. I think 6 is more appropriate, comparing to the ZPU-4 battery.

        \begin{notesources}
            \item Wikipedia on the \href{https://en.wikipedia.org/wiki/ZU-23-2}{ZU-23}.
            \item WeaponSystems on the \href{https://weaponsystems.net/system/116-ZU-23-2}{ZU-23} and \href{https://weaponsystems.net/system/1580-23mm\%202A14}{2A14}.
        \end{notesources}
    }

    \subsectionwithcounter{mobile ZU-23 section}{Mobile ZU-23 Section}

    This is a section of two light trucks mounting ZU-23 guns.
    As with the earlier ZPU-1/2/4, many users of the ZU-23 mount it on vehicles, including trucks, tracked carriers, and APCs, for better mobility.
    Informal conversions are typical used as \emph{technicals} in ground combat.%
    \note{Mobile ZU-23 Section}{

        \aaamountvalues{23}{2000}{980}

        \input{aaavalues/mobile ZU-23 section.tex}

        I adopt the AAA values, defense strength, and sighting range from {\AirStr} and {\MPC}, but the hit rolls and damage are taken from the ZU-23 battery with a \minus{1} modifier to the hit rolls for being a section.

        The hit rolls and damage in {\AirStr} and {\MPC} are 2/2/1 and 2. As with the ZU-23 battery, these seem low. The M16 has a similar rate of fire and smaller caliber, but much better hit rolls and the same damage rating, and this does not seem appropriate. Therefore, by analogy with the M16 and in agreement with the empirical model, I give it hit rolls of 4/3/2 and a damage rating of 3.

        %The VPs seem about right, as the battery has 6 VPs. {\AirStr} gives 5 VPs.

        \begin{notesources}
            \item Wikipedia on the \href{https://en.wikipedia.org/wiki/ZU-23-2}{ZU-23} and \href{https://en.wikipedia.org/wiki/Technical_(vehicle)}{Technicals}.
            \item WeaponSystems on the \href{https://weaponsystems.net/system/116-ZU-23-2}{ZU-23} and \href{https://weaponsystems.net/system/1580-23mm\%202A14}{2A14}.
        \end{notesources}
    }

    \subsectionwithcounter{BTR-152A ZU-23 section}{BTR-152A ZU-23 Section}

    The BTR-152A ZU-23 has a twin 23~mm ZU-23 turret in the open, six-wheeled BTR-152 APC.
    A section has two vehicles.
    It entered service with the PLO prior to the 1982 War in Lebanon; the in-service data is a guess.%
    \note{BTR-152A ZU-23 Section}{

        \aaamountvalues{23}{2000}{980}

        \input{aaavalues/BTR-152A ZU-23 section.tex}

        I adopt the AAA properties of the mobile ZPU-2/4 section and the defense strength and sighting range of the M16.
        The adopted AAA values are in agreement with the empirical model.

        The BTR-152 is based on a truck, and so I give it “G/P” mobility rather than “U”.

        Army Guide states that the A/E mounted a ZPU-2 and the D a ZPU-4. Wikipedia states that the A/D/E referred to the chassis (with the D using the BTR-152V and the E the BTR-152V1 chassis) rather than the gun mount. I am not sure how to resolve this, but I don't think it is that important. Instead, I will use “BTR-152A ZPU-2” and “BTR-152A ZPU-4,” which extends straightforwardly to other guns mounted on the BTR-152.

        \begin{notesources}
            \item Wikipedia on the \href{https://en.wikipedia.org/wiki/BTR-152\#Soviet_Union}{BTR-152A} and \href{https://en.wikipedia.org/wiki/ZU-23-2}{ZU-23}.
            \item Army Guide on the \href{https://www.army-guide.com/eng/product1120.html}{BTR-152}.
            \item Miranda, M., 2015: “\href{https://thesovietarmourblog.blogspot.com/2015/12/}{Gun Truck}”
        \end{notesources}
    }

    \subsectionwithcounter{ZSU-23-4 section}{ZSU-23-4 Section}

    The ZSU-23-4 has four radar-controlled 23~mm guns, similar to those in the ZU-23, mounted in an armored turret on a hull adapted from the PT-76 light tank.
    A section has two vehicles.
    The ZSU-23-4 entered service in the Soviet Army in 1965, replacing the ZSU-57-2 and outclassing all NATO anti-aircraft guns at that time. A platoon of four was assigned to the air-defense battery of tank and motor-rifle regiments and later complemented by a platoon of four SA-9 vehicles. It was exported to Soviet allies and saw extensive combat with Arab forces in the 1973 Arab-Israeli War and with Iraqi forces in the Iran-Iraq War and the 1991 Gulf War.%
    \note{ZSU-23-4 Section}{

        \aaamountvalues{23}{4000}{980}

        \input{aaavalues/ZSU-23-4 section.tex}

        I adopt the AAA values, defense strength, and sighting range from {\AirStr} and {\MPC}.
        The adopted AAA values are in agreement with the empirical model.

        RadarTutorial states that the radar operates at 14.6--15.6 GHz (NATO J-band and IEEE Ku-band), and does not mention a difference between search and tracking. This makes sense, as there seems to be only one radar dish.

        Wikipedia states that the ZSU-23-4M3 “Biryusa” upgrade from 1977 adds IFF.

        \begin{notesources}
            \item Wikipedia on the \href{https://en.wikipedia.org/wiki/ZSU-23-4_Shilka}{ZSU-23-4}.
            \item WeaponSystems on the \href{https://weaponsystems.net/system/452-ZSU-23-4\%20Shilka}{ZSU-23-4} and \href{https://weaponsystems.net/system/1580-23mm\%202A14}{2A7}.
            \item RadarTutorial on the \href{https://www.radartutorial.eu/19.kartei/11.ancient4/karte088.en.html}{RPK-2/1RL33} radar.
        \end{notesources}
    }

    \subsectionwithtwocounters{Sinai 23 Ayn-al-Saqr section}{Sinai 23 Stinger section}{Sinai 23 Section}

    The Sinai 23 has a Dassault TA20 turret with two 23~mm guns and six Ayn-al-Saqr (similar to the Strela-2M) or Stinger missiles on an M113 chassis.
    A section has two vehicles.
    Each battery of two sections operates with one vehicle equipped with an RA~20S search and ranging radar.
    The Sinai 23 entered service with the Egyptian Army in 1992.%
    \note{Sinai 23 Section}{

        \aaamountvalues{23}{2000}{980}

        Sinai 23 Ayn-al-Saqr \input{aaavalues/Sinai 23 Ayn-al-Saqr section.tex}
        Sinai 23 Stinger \input{aaavalues/Sinai 23 Stinger section.tex}

        I adopt the AAA properties of mobile ZU-23, but with \plus{1} to the hit rolls for the FCR-R giving 5/4/3. This means that the Sinai 23 (with two barrels) has the same hit rolls and damage rating as the ZSU-23-4 (with four barrels). Perhaps this is acceptable, since there are 30 years between these systems.
        {\MP} gives this hit rolls of 3/3/2 and no FCR-R. With FCR-R, these hit rolls would be 4/4/3, which is close to those adopted.
        The adopted AAA values are in agreement with the empirical model.
        I adopt the defense strength and sighting range from {\EOTG} and {\MPC}.

        The turret is the Dassault TA-23E and seems to be similar to the LAV-AD turret, but with a pair of 23~mm guns instead of a GAU-12 25~mm gun.

        I don't know which model of Stinger was used; I would guess the FIM-92A.

        {\MP} has the Sinai 23 with SA7B from 1990 and the Nile 23 with Stingers from 1998. Wikipedia states that they were similar but competing vehicles, and the Sinai 23 was selected in preference to the Nile 23.

        The brochure explicitly mentions a “day sight”; there is no sign of an IR sight.

        A platoon seems to consist of three “satellite” fire units and one “leader” unit with a RA 20S radar. The radar has only 1 degree precision, so it is not a traditional FCR, but seems to give warning, cueing, and range for optical tracking. I will treat it as an FCR-R and give a \plus{1} modifier.

        The brochure states radar works in the E-band and is pulse-Doppler. There is no mention of IFF.

        Perhaps the most accurate way to handle the radar would be to have a platoon of 3 fire units and 1 radar unit, and for each D determine if the radar unit is killed on a \minusafter{3}.

        \begin{notesources}
            \item Wikipedia on \href{https://en.wikipedia.org/wiki/ZU-23-2\#Egypt}{Sinai 23} and \href{https://en.wikipedia.org/wiki/ZU-23-2}{ZU-23}.
            \item \href{https://forum.warthunder.com/t/british-weapon-systems-technical-data-and-discussion/3879/2233}{Brochure} describing the Sinai 23.
            \item \href{https://forum.warthunder.com/t/british-weapon-systems-technical-data-and-discussion/3879/2233}{Description} of RA 20S FCR.

        \end{notesources}
    }
    \vspacelines{1}

    \subsectionwithcounter{LAV-AD section}{LAV-AD Section}

    The LAV-AD has a 25~mm GAU-12 rotary cannon and eight FIM-92D Stinger missiles mounted in a turret on the LAV-25 vehicle.
    A section has two vehicles.
    The LAV-AD does not have radar, but is equipped with a passive night IR sight.
    It entered service with the USMC in 1997, but was retired after only a few years.%
    \note{LAV-AD Section}{

        \aaamountvalues{25}{1800}{1040}

        \input{aaavalues/LAV-AD section.tex}

        I have adopted the AAA value, defense strength, and sighting range of the Blazer from {\AirStr}, but without FCR-W (\minus{2} to the hit rolls), with a laser range finder (\plus{1} to the hit rolls), and with night IR sights.
        {\MP} gives the same values, but reduces the maximum altitude from 10 to 9; I have not adopted this change.
        The adopted AAA values are in agreement with the empirical model.

        \begin{notesources}
            \item Wikipedia on the \href{https://en.wikipedia.org/wiki/LAV-25\#Other_Variants}{LAV-AD}, \href{https://en.wikipedia.org/wiki/GAU-12_Equalizer}{GAU-12}, and \href{https://en.wikipedia.org/wiki/FIM-92_Stinger}{FIM-92}.
            \item Dickerson, C.R., 1991: “\href{https://www.globalsecurity.org/military/library/report/1991/DCR.htm}{LAV-AD: Mobile Air Defense for the Ground Maneuver Force}”.
        \end{notesources}
    }

    \subsectionwithcounter{M5359 section}{M53/59 Section}

    The M53/59 has a twin M53 30~mm gun mount on a partially armored truck. The fire-control system is entirely optical. A section has two vehicles.
    It  entered service in the Czech Army in 1959 and was used instead of the ZSU-57-2. It was also exported to Iraq and saw combat in the Iran-Iraq War and the 1991 Gulf War.%
    \note{M53/59 Section}{

        \aaamountvalues{30}{1000}{1000}

        \input{aaavalues/M5359 section.tex}

        I adopt the AAA values from {\AirStr} and {\MPC}, but reduce the hit rolls from 4/3/2 to 3/2/1.
        I adopt the defense strength (4), target class (open), and sighting range (12) from the M163.

        The hit rolls for the M53/59 section from {\AirStr} and {\MPC} are 4/3/2. The empirical model suggests 3/2/1, derived from the modified hit rolls of 4/3/2 for the Oerlikon GDF-001 battery with a \minus{1} modifier for being a section and also from the hit rolls of 5/3/2 for the AMX-30SA with a \minus{2} modifier to discounting the FCR-W. Therefore, by analogy with these units and in agreement with the empirical model, I adopt hit rolls of 3/2/1.

        The M53/59 is based on a truck, and so I give it “G/P” mobility rather than “U”.

        \begin{notesources}
            \item Wikipedia on the \href{https://en.wikipedia.org/wiki/M53/59_Praga}{M53/59}.
            \item WeaponSystems on the \href{https://weaponsystems.net/system/751-EE03\%20-\%20M53/59}{M53/59} and \href{https://weaponsystems.net/system/752-M53}{M53}.
        \end{notesources}
    }

    \subsectionwithcounter{AMX-13 DCA section}{AMX-13 DCA Section}

    The AMX-13 DCA is equipped with a pair of radar-controlled 30~mm Hispano-Suiza HS.831A guns mounted in an armored turret on the hull of an AMX-13 tank.
    A section has two vehicles.
    The AMX-13 DCA is equipped with the Thomson-CSF Oeil Noir radar.
    It served only in the French Army, entering service in 1969 and being retired in the 1990s.%
    \note{AMX-13 DCA Section}{

        \aaamountvalues{30}{1200}{970}

        \aaavalues{11}{4/7/9}{4/3/2}{3}{R/VF}


        I adopt the {\AirStr} values given for the AMX-30 DCA, which has identical guns.

        {\MPC} gives slightly different values for different 30 mm systems. For the Tunguska and M53/59, it gives a maximum altitude of 11 and ranges of 4/7/9. For the AMX-13, BOV 30, and K-30 it gives 10 and 3/6/9. However, there seems to be very little difference between the muzzle velocities and masses of the shells fired by these guns. I prefer to keep the {\AirStr} values for the AMX-13 DCA, which match those he gives for the Tunguska and M53/59.

        {\MPC} also gives a damage rating of 4. This is the same as the Tunguska, despite the latter having a much higher rate of fire. I prefer to use the {\AirStr} value for the AMX-30 DCA of 3.

        I believe the CGT/Thomson-CSF \foreignlanguage{french}{Œil Noir} DRVC 1A radar provides search and ranging, but not tracking. BAS'ART states:
        \begin{quote}
            \foreignlanguage{french}{Distance fournie par le radar en mode télémétrie. Vitesse de rotation de la droite pièce-objectif fournie par le pointage sur l’objectif.}
        \end{quote}
        I translate this as:
        \begin{quote}
            Distance furnished by the radar in telemetry mode.
            Tracking speed of the target furnished by the pointing of the sight.
        \end{quote}
        This is consistent with the radar of the Panhard M3 VDA, which is more or less contemporaneous, and with later statements that the AMX-30 DCA did not have tracking but the AMX-30 SA did.

        BAS'ART states that the radar operates at 1.71--1.75 GHz (NATO D-band and IEEE S-band) and is a PD radar with MTI.

        \begin{notesources}
            \item \href{https://artillerie.asso.fr/basart/article.php3?id_article=1086}{BAS'ART} on the AMX-13 “\foreignlanguage{french}{Bitubes de 30 mm automoteur antiaérien}”.
            \item Wikipedia on the \href{https://en.wikipedia.org/wiki/AMX-13\#Production_variants}{AMX-13 DCA} and \href{https://en.wikipedia.org/wiki/Oerlikon_KCB}{HS.831A}.
            \item WeaponSystems on the \href{https://weaponsystems.net/system/114-AMX-30+DCA}{AMX-13 DCA} \href{https://weaponsystems.net/system/1562-30mm\%20Hispano-Suiza\%20HS.831}{HS.831A}.
        \end{notesources}
    }

    \subsectionwithcounter{AMX-30SA section}{AMX-30SA Section}

    The AMX-30SA is equipped with a pair of radar-controlled 30~mm Hispano-Suiza guns mounted in an armored turret on the hull of an AMX-30 tank.
    A section has two vehicles.
    AMX-30SA is an improved version of the AMX-30 DCA and served only in the Saudi Army from 1979.%
    \note{AMX-30SA Section}{

        \aaamountvalues{30}{1200}{970}

        \aaavalues{11}{4/7/9}{5/4/3}{3}{R/VF}


        I adopt the {\AirStr} values given for the AMX-30 DCA. The real DCA does not have a tracking radar, but the {\AirStr} one does, matching the SA.

        \begin{notesources}
            \item Wikipedia on the \href{https://en.wikipedia.org/wiki/AMX-30\#AMX-30_DCA_–_défense_contre_avion}{AMX-30DCA/SA} and \href{https://en.wikipedia.org/wiki/Oerlikon_KCB}{HS.831A}.
            \item WeaponSystems on the \href{https://weaponsystems.net/system/114-AMX-30+DCA}{AMX-30DCA} \href{https://weaponsystems.net/system/1562-30mm\%20Hispano-Suiza\%20HS.831}{HS.831A}.
        \end{notesources}
    }

    \subsectionwithcounter{Tunguska section}{Tunguska Section}

    The Tunguska has twin radar-controlled 30mm guns and four SA-19 missiles in an armored turret on an armored and tracked hull.
    A section has two vehicles.
    The Tunguska was designed to counter the A-10 and AH-64, against which the earlier ZSU-23-4 had weaknesses.
    It entered service in the Soviet Army in 1984. It has been used by Russia, Ukraine, and Belarus after the break-up of the Soviet Union and has been exported.
    The Tunguska is referred to as the “ZSU-30-2” in {\AirStr}.%
    \note{Tunguska Section}{

        \aaamountvalues{30}{5000}{960}

        \aaavalues{11}{4/7/9}{5/4/3}{4}{W/VF}


        I adopt the AAA values from {\AirStr} and {\MP}. The Tunguska is called the ZSU-30-2 in {\AirStr}.

        %The VP values are tentative, and might be modified once other SAM launchers are considered.

        Wikipedia states that 1RL144 radar system uses the E-band for search, the J-band for tracking, and has IFF. RadarTutorial states it uses S and X band.

        \begin{notesources}
            \item Wikipedia on the \href{https://en.wikipedia.org/wiki/2K22_Tunguska}{Tunguska}.
            \item WeaponSystems on the \href{https://weaponsystems.net/system/60-2S6\%20Tunguska}{Tunguska}, \href{https://weaponsystems.net/system/558-30mm+2A38}{2A38}, and \href{https://weaponsystems.net/system/1433-9M311}{9M311}.
            \item RadarTutorial on the \href{https://www.radartutorial.eu/19.kartei/04.battle/karte057.en.html}{1RL144}.
        \end{notesources}
    }

    \subsectionwithcounter{Tunguska-M section}{Tunguska-M Section}

    The Tunguska-M is an improved version of the Tunguska with eight missiles instead of four. A section has two vehicles. It entered service in the Soviet Army in 1990. It has been used by Russia, Ukraine, and Belarus after the break-up of the Soviet Union and has been exported.%
    \note{Tunguska-M Section}{

        \aaamountvalues{30}{5000}{960}

        \input{aaavalues/Tunguska-M section.tex}

        I adopt the AAA values from {\AirStr} and {\MP}. The Tunguska is called the ZSU-30-2 in {\AirStr}.
        %The VP values are tentative, and might be modified once other SAM launchers are considered.

        \begin{notesources}
            \item Wikipedia on the \href{https://en.wikipedia.org/wiki/2K22_Tunguska}{Tunguska}.
            \item WeaponSystems on the \href{https://weaponsystems.net/system/60-2S6\%20Tunguska}{Tunguska}, \href{https://weaponsystems.net/system/558-30mm+2A38}{2A38}, and \href{https://weaponsystems.net/system/1433-9M311}{9M311}.
        \end{notesources}
    }

    %\paragraph{Pantsir..} The Pantsir has twin radar-controlled 30~mm guns and twelve SA-22 missiles in an unarmored turret mounted on a truck. While the Tunguska was designed to accompany ground forces, the Pantsir was designed for point defense of mobile S-300/S-400 batteries and other high-value targets, and this led to the choice of the system being wheeled and unarmored. It entered service in 2012 with Russian forces and has since been exported.

\end{withcounters}

\section{Medium AAA Units}

\begin{withcounters}

    \subsectionwithcounter{Oerlikon GDF battery}{Oerlikon GDF Battery}

    The Oerlikon GDF has two 35~mm guns in an open mount on a towed carriage.
    A battery has six gun mounts.
    The basic fire-control system is optical, but a section is often coupled with the Super Fledermaus or Skyguard fire-control radars.
    (These radars cannot control a whole battery, only a section.)
    The Oerlikon GDF entered service in 1963 and has served widely. Notably, it saw combat in the South Atlantic War with the Argentine forces in the defense of BAM Malvinas (Port Stanley Airport) and BAM Condor (Goose Green Airfield).%
    \note{Oerlikon GDF Battery}{

        \aaamountvalues{35}{1100}{1175}

        \input{aaavalues/Oerlikon GDF battery.tex}

        The AAA values for the battery are from {\AirStr} with the hit rolls reduced from 5/4/3 to 4/3/2 as described in the extended note above.

        \begin{notesources}
            \item Wikipedia on the \href{https://en.wikipedia.org/wiki/Oerlikon_GDF}{GDF}.
            \item WeaponSystems on the \href{https://weaponsystems.net/system/933-EE02\%20-\%20GDF}{GDF} and \href{https://weaponsystems.net/system/283-35mm\%20Oerlikon\%20KDA}{KDA}.
        \end{notesources}
    }

    \subsectionwithcounter{Gepard section}{Gepard Section}

    The Gepard has two radar-controlled 35~mm guns in an armored turret on an adapted Leopard 1 tank hull. A section has two vehicles.
    The Gepard entered service in 1976 with the West German Army and also served with the Belgian and Dutch armies. After the end of the Cold War, they were retired by these armies, and many were sold on to other countries.%
    \note{Gepard Section}{

        \aaamountvalues{35}{2200}{1175}

        \input{aaavalues/Gepard section.tex}

        The AAA values are from {\AirStr}.

        Wikipedia states that the Siemens MPRD-12 search radar operates in the S band (2.60--3.95 GHz) and the tracking radar in the Ku band (12.4--18.0 GHz). RadarTutorial states that it searches in the E band (2--3 GHz) and tracks in the J band (10--20 GHz) and has an integrated IFF.

        \begin{notesources}
            \item Wikipedia on the \href{https://en.wikipedia.org/wiki/Flakpanzer_Gepard}{Gepard}.
            \item WeaponSystems on the \href{https://weaponsystems.net/system/69-Gepard}{Gepard} and \href{https://weaponsystems.net/system/283-35mm\%20Oerlikon\%20KDA}{KDA}.
            \item RadarTutorial on the \href{https://www.radartutorial.eu/19.kartei/11.ancient/karte031.en.html}{MPRD-12}.
        \end{notesources}
    }

    \subsectionwithcounter{61-K battery}{61-K Battery}

    The 61-K (M1939) has a single 37~mm gun in an open mount on a towed carriage.
    A battery has six gun mounts.
    The fire-control system is entirely optical, and (as an exception to the normal rule allowing medium AAA units to employ external FCRs) the 61-K cannot be coupled to an external fire-control radar.
    After entering service in 1939 with the Soviet Army, it served in WW2 and after until replaced by the S-60 57~mm gun. It also served with many Soviet allies, and in particular with communist forces in the Korean and Vietnam Wars.
    The 61-K is referred to as the “M-38” in {\AirStr}, {\EOTG}, and {\TSOH}.%
    \note{61-K Battery}{

        \aaamountvalues{37}{165}{880}

        \aaavalues{12}{4/8/11}{3/2/1}{4}{}


        The AAA values are from {\AirStr} and {\TSOH}.
        % but TSOH gives it 6 VPs and {\AirStr} only 5.

        I can find no sources other than {\AirStr}, {\EOTG}, and {\TSOH} that refer to the 61-K 37~mm gun as the “M-38”. Another designation for the gun is “M-1939”, so I could understand “M-39”.

        \begin{notesources}
            \item Wikipedia on the \href{https://en.wikipedia.org/wiki/37_mm_automatic_air_defense_gun_M1939_(61-K)}{61-K}.
            \item WeaponSystems on the \href{https://weaponsystems.net/system/950-M1939}{61-K}.
        \end{notesources}
    }

    \subsectionwithcounter{M15 section}{M15 Section}

    The M15 Combination Gun Motor Carriage has a M1 37-mm cannon and two Browning M2 .50 cal (12.7 mm) machine guns in an M42 turret mounted on a modified M3 half-track.
    A section has two vehicles.
    The M15 entered service with the US Army in 1942, saw combat in WW2 and the Korean War, in the latter mainly in a ground-support role, and was retired in 1953. It also served with the armed forces of Japan and Yugoslavia.%
    \note{M15 Section}{

        \aaamountvalues{37}{120}{792}

        \aaavalues{12}{2/4/8}{3/2/1}{4}{}


        The M1 is somewhat inferior to the 61-K in rate of fire and muzzle velocity. There is a quote on Wikipedia about using the .50 cal machine guns to aim the 37 mm cannon.

        Considering this as half an M16, the ranges would be 2/3/4 and the hit rolls 3/2/1. Considering this as half a 61-K battery, the ranges would be 4/8/11 and the hit rolls 2/1/0. We can combine these two by giving it ranges of 2/4/8 and hit rolls of 3/2/1. The damage rating is 4 to match the 37 mm.

        \begin{notesources}
            \item Wikipedia on the \href{https://en.wikipedia.org/wiki/M15_half-track}{M15}, \href{https://en.wikipedia.org/wiki/37_mm_gun_M1}{M1 37 mm}, and \href{https://en.wikipedia.org/wiki/M2_Browning}{M2 .50 cal}.
        \end{notesources}
    }

    \subsectionwithcounter{Bofors L60 battery}{Bofors L/60 Battery}

    The Bofors L/60 has a single 40~mm gun in an open mount on a towed carriage. A
    battery has six gun mounts.
    The fire-control system is entirely optical.
    The Bofors L/60 entered service in the 1940s and was widely used by Allied forces during WW2, both on land and sea. After the war, many continued to see service, although it was increasingly inadequate against faster jet aircraft and in many cases was replaced by the new L/70 model.%
    \note{Bofors L/60 Battery}{

        \aaamountvalues{40}{135}{865}

        \aaavalues{15}{4/8/12}{2/2/1}{4}{}


        The AAA values are from {\EOTG}.

        \begin{notesources}
            \item Wikipedia on the \href{https://en.wikipedia.org/wiki/Bofors_40_mm_L/60_gun}{Bofors L/60}.
            \item WeaponSystems on the \href{https://weaponsystems.net/system/1574-Bofors\%20L/60}{Bofors L/60}.
        \end{notesources}
    }

    \subsectionwithcounter{M19 section}{M19 Section}

    The M19 Multiple Gun Motor Carriage has two Bofors 40~mm L/60 guns in a partially armored mount on a tracked chassis based on that of the M24 light tank.
    A section has two vehicles.
    The fire-control system is entirely optical.
    The M19 entered service in 1944 with the US Army, but did not see combat in WW2. It was used in the Korean War, mostly in against ground targets, but was retired in 1953 and replaced by the M42. It served in the Netherlands Army from 1951 to 1978.%
    \note{M19 Section}{

        \aaamountvalues{40}{270}{865}

        \input{aaavalues/M19 section.tex}

        The AAA values are identical to those of the Bofors L/60 battery, since the number of guns is the same.

        \begin{notesources}
            \item Wikipedia on the \href{https://en.wikipedia.org/wiki/M19_Multiple_Gun_Motor_Carriage}{M19}.
        \end{notesources}
    }

    \subsectionwithcounter{M34 section}{M34 Section}

    The M34 Gun Motor Carriage has a Bofors 40 mm L/60 cannon in an open mount on a modified M15 half-track.
    Many M15s were converted into M34s as supplies of 37-mm shells for the M1 cannon became scarce. Unlike the M15, the gun crew have no armored protection.
    A section has two vehicles.
    The M34 entered service with the US Army in 1951, saw combat in the Korean War, mainly in a ground-support role, and was likely retired in 1953.%
    \note{M34 Section}{

        \aaamountvalues{40}{135}{865}

        \aaavalues{15}{2/4/8}{1/1/1}{4}{}


        Considering this as half an M19, the ranges would be 4/8/12 and the hit rolls 1/1/0. Since no FCR can be added, the effective range is 8. We therefore give it ranges of 2/4/8 and hit rolls of 1/1/1. The damage rating is 4 to match the M19.

        \begin{notesources}
            \item Wikipedia on the \href{https://en.wikipedia.org/wiki/M15_half-track}{M34} and \href{https://en.wikipedia.org/wiki/Bofors_40_mm_L/60_gun}{Bofors 40 mm L/60}.
            \item AFV Data Base on the \href{https://afvdatabase.com/usa/40mmgmcm34.html}{M34}.
        \end{notesources}
    }

    \subsectionwithcounter{Bofors L70 battery}{Bofors L/70 Battery}

    The Bofors L/70 has a single 40~mm gun in an open mount on a towed carriage.
    A battery has six gun mounts.
    Between the 1930s and the end of WW2, aircraft speed increased substantially, and so the engagement time of the Bofors L/60 became inadequate. The new L/70 fired a lighter shell at a higher velocity and achieved a significant improvement in range and engagement time. The basic fire-control system is optical, but various external fire-control radars can be used with it, including the Super Fledermaus and Flycatcher systems.
    NATO forces and others adopted it widely starting in 1952. Israel used batteries with the Super Fledermaus FCR for point-defense in the 1967 war.%
    \note{Bofors L/70 Battery}{

        \aaamountvalues{40}{330}{1000}

        \aaavalues{15}{4/8/12}{3/2/1}{4}{}


        The AAA values are from {\AirStr}.

        Singh (ch. 4) describes the Israeli use.

        % LTeX: dictionary+=Defence
        \begin{notesources}
            \item Wikipedia on the \href{https://en.wikipedia.org/wiki/Bofors_40_mm_Automatic_Gun_L/70}{Bofors L/70}.
            \item WeaponSystems on the \href{https://weaponsystems.net/system/932-Bofors\%20L/70}{Bofors L/70}.
            \item Singh, M. (Pen \& Sword, 2020): “Anti-Aircraft Artillery in Combat 1950--1972: Air Defence in the Jet Age”.
        \end{notesources}
        % LTeX: dictionary-=Defence

    }

    \subsectionwithcounter{M42 section}{M42 Section}

    The M42 Duster has two Bofors 40~mm L/60 guns in a partially armored mount on a tracked chassis based on that of the M41 light tank.
    A section has two vehicles.
    The fire-control system is entirely optical.
    The M42 entered service in 1953 with the US Army and US Marine Corps, replacing the M19, saw combat in the Vietnam War against ground targets, and served at least with reserve units until 1988. It was also used by Austria, Greece, Japan, Jordan, Lebanon, Pakistan, Taiwan, Thailand, Tunisia, Turkey, Venezuela, Vietnam (South and North), and West Germany.%
    \note{M42 Section}{

        \aaamountvalues{40}{270}{865}

        \input{aaavalues/M42 section.tex}

        The AAA values are identical to those of the Bofors L/60 battery, since the number of guns is the same.

        \begin{notesources}
            \item Wikipedia on the \href{https://en.wikipedia.org/wiki/M42_Duster}{M42}.
        \end{notesources}
    }

    \subsectionwithcounter{Bofors L70 BOFI battery}{Bofors L/70 BOFI}

    The BOFI version of the Bofors L/70 adds modern sights, a laser rangefinder, and proximity-fuzed shells. A battery has six gun mounts. The BOFI is used by the armed forces of Brazil.%
    \note{Bofors L/70 BOFI Battery}{

        \aaamountvalues{40}{330}{1000}

        \aaavalues{15}{4/8/12}{5/5/4}{4}{}


        The AAA values are adapted from {\AirStr} values for the Bofors L/70, with the hit rolls increased by one for proximity fuzes and one for the BOFI laser rangefinder (effectively FCR-R), giving hit rolls of 5/5/4 for BOFI.

        I don't know when the BOFI became available or who else uses them.

        \begin{notesources}
            \item Wikipedia on the \href{https://en.wikipedia.org/wiki/Bofors_40_mm_Automatic_Gun_L/70\#Variants}{BOFI and BOFI-R}.
            \item WeaponSystems on the \href{https://weaponsystems.net/system/932-Bofors\%20L/70}{BOFI and BOFI-R}.
        \end{notesources}
    }

    \subsectionwithcounter{Bofors L70 BOFI-R battery}{Bofors L/70 BOFI-R Battery}

    The BOFI-R version of the Bofors L/70 also adds an integrated fire-control radar to the BOFI version. A battery has six gun mounts.%
    \note{Bofors L/70 BOFI-R Battery}{

        \aaamountvalues{40}{330}{1000}

        \aaavalues{15}{4/8/12}{6/5/4}{4}{W/VF}


        The AAA values are adapted from {\AirStr} values for the Bofors L/70, with the hit rolls increased by one for proximity fuzes and another two for the BOFI-R radar, giving hit rolls of 6/6/5 for BOFI-R whereas {\AirStr} gives the BOFI-R 5/5/4.

        I don't know when the BOFI-R became available or who else uses them.

        Wikipedia states the BOFI-R has: “Multisensor fire control system with a J band radar. Provides automatic acquisition and tracking with an effective range of 4 km without external radar input.”

        \begin{notesources}
            \item Wikipedia on the \href{https://en.wikipedia.org/wiki/Bofors_40_mm_Automatic_Gun_L/70\#Variants}{BOFI and BOFI-R}.
            \item WeaponSystems on the \href{https://weaponsystems.net/system/932-Bofors\%20L/70}{BOFI and BOFI-R}.
        \end{notesources}
    }

    \subsectionwithcounter{S-60 battery}{S-60 Battery}

    The S-60 is a Soviet single 57~mm gun on an open mount on a towed carriage.
    A battery has six gun mounts.
    The basic fire-control system is optical, but the gun is often used with the SON-9 (Fire Can) radar.
    The S-60 entered service in 1950 with Soviet forces, replacing the 61-K 73~mm gun.
    It was exported to many Soviet allies and saw combat with Arab forces in the 1967 and 1973 Arab-Israeli War, with communist forces in the Vietnam War, and with Iraqi forces in the Iran-Iraq War and the Gulf War. However, it was not used by communist forces in the Korean War. When paired with the SON-9 (Fire Can) radar, it was considered one of the more dangerous AAA systems faced by US aircraft in Vietnam.%
    \note{S-60 Battery}{

        \aaamountvalues{57}{110}{1000}

        \aaavalues{18}{5/10/15}{2/2/1}{5}{}


        The AAA values are from {\AirStr} and {\TSOH}.
        % These are the same as in {\AirStr}, except that TSOH gives 8 VPs and {\AirStr} only 6.

        \begin{notesources}
            \item Wikipedia on the \href{https://en.wikipedia.org/wiki/AZP_S-60}{S-60}.
            \item WeaponSystems on the \href{https://weaponsystems.net/system/477-S-60}{S-60}.
        \end{notesources}
    }

    \subsectionwithcounter{ZSU-57-2 section}{ZSU-57-2 Section}

    The ZSU-57-2 has two 57~mm cannons in a lightly armored but open-topped turret on an armored and tracked hull.
    A section has two vehicles.
    The cannons are similar to the S-60 cannons and the fire-control system is optical.
    The ZSU-57-2 entered service in the Soviet Army in 1955 and replaced the BTR-40A and BTR-152A self-propelled 14.5~mm AA guns in the AA batteries of tank regiments. It was in turn replaced by the ZSU-23-4 from 1965. The ZSU-57-2 also served widely with Soviet allies and saw combat with Egypt and Syria in the 1967 and 1973 Wars, with Syria in the 1982 Lebanon War, with North Vietnam in the 1972 Easter Offensive and the 1975 Spring Offensive, and with various factions in the Yugoslav Wars in the 1990s.%
    \note{ZSU-57-2 Section}{

        \aaamountvalues{57}{240}{1000}

        \aaavalues{18}{5/10/15}{2/1/1}{5}{}


        For the ZSU-57-2 section: The AAA values are from {\EOTG}.

        \begin{notesources}
            \item Wikipedia on the \href{https://en.wikipedia.org/wiki/ZSU-57-2}{ZSU-57-2}.
            \item WeaponSystems on the \href{https://weaponsystems.net/system/531-ZSU-57-2\%20Sparka}{ZSU-57-2}.
        \end{notesources}
    }

    \subsectionwithcounter{M51 battery}{M51 Battery}

    The M51 Skysweeper has a single 75~mm gun in an open mount in a towed carriage.
    A battery has four mounts.
    The mount incorporates the T38 radar fire-control system and the gun has a relatively high rate of fire of 45 rounds per minute.
    Proximity-fuzed shells are available.
    The M51 entered service with the US Army in 1951 and was used for defense of strategic locations in the USA and in Europe, but was retired in 1959. It also served with Greece, Japan, and Turkey, with the last examples being retired until the mid-1970s.
    \note{M51 Battery}{

        \aaamountvalues{75}{45}{854}

        \aaavalues{30}{6/12/18}{4/4/3}{5}{W/MF}


        Wikipedia gives a rate of fire of 45 rounds per minute, a muzzle velocity of 854~m/s, an effective maximum altitude of 30k feet, and a horizontal range of 26 hexes (13 km).

        The TM states that the search radar is X-band (6--11 GHz). RadarTutorial also states it is X-band or I/J-band. There is no mention of PD or MTI.

        RadarTutorial notes that it is often used with the AN/MPQ-10 as a search radar, but this seems to be counter-battery radar.

        The AAA hit rolls are based on the S-60, reduced by 1 for a lower rate of fire, increased by 1 for proximity fuzes, and increased by 2 for FCR-W, so 4/4/3. I will give it the same range of 6/12/18 as the 85~mm JS-12, which fires a slightly heavier round at a slightly lower velocity.

        \begin{notesources}
            \item Wikipedia on the \href{https://en.wikipedia.org/wiki/M51_Skysweeper}{M51} and \href{https://en.wikipedia.org/wiki/Western_Electric_M-33_Antiaircraft_Fire_Control_System}{T38 radar}
            \item WeaponSystems on the \href{https://weaponsystems.net/system/421-M51+Skysweeper}{M51}.
            \item Werrell, K.P. (2005, Air University Press, 2nd edition): “\href{https://media.defense.gov/2017/Mar/31/2001725225/-1/-1/0/B_0028_WERRELL_ARCHIE_TO_SAM.PDF}{Archie to SAM}”.
            \item Department of the Army (1956): \href{https://www.usarmygermany.com/Sont.htm?http&&&www.usarmygermany.com/Units/Air\%20Defense/USAREUR_ADA\%20Overview\%201.htm}{TM 9-6092-1-1} “Antiaircraft Fire Control System M33C and M33D Operation, June 1956”.
            \item RadarTutorial on the \href{https://www.radartutorial.eu/19.kartei/11.ancient7/karte098.en.html}{T-38} and \href{https://www.radartutorial.eu/19.kartei/11.ancient3/karte037.en.html}{AN/MPQ-10}.
        \end{notesources}
    }

\end{withcounters}

\section{Heavy AAA Units}

\begin{withcounters}

    \subsectionwithcounter{Flab Kan 38 battery}{Flab Kan 38 Battery}

    The Swiss Flab Kan 38 (Fliegerabwehr Kanone 38) has a single 75~mm L/49 gun in an open mount on a towed carriage.
    It is a licensed copy of the French Canon de 75 contre aéronefs modèle 1939 Schneider.
    A battery has six gun mounts.
    Its fire control system is optical, but in Swiss service from 1958 it was paired with the Mk.7 Blue Cedar radar.
    The Flab Kan 38 entered service with the Swiss Air Force in 1939 and was retired in 1967.%
    \note{Flab Kan 38 Battery}{

        \aaamountvalues{75}{25}{700}

        \aaavalues{24}{6/12/18}{2/1/0}{5}{}


        Wikipedia gives the 75 mm Schneider a rate of fire of 20--30 RPM and a muzzle velocity of about 700 m/s.

        Wikipedia gives an effective altitude of 27,000 ft. However, the KS-12 fires a heavier shell at a higher velocity and has a maximum altitude of 27,000 ft. To match this, I will give it 24,000 ft.

        The AAA ranges match those of the KS-12. I presume it does not have proximity fuzes, so I will give it hit rolls of 2/1/0.

        I classify the Flab Kan 38 as heavy but the M51 as medium, despite both having the same caliber. This is because the M51 has powered traverse whereas the Flab Kan 38 has manual traverse.

        \begin{notesources}
            \item Wikipedia on the \href{https://en.wikipedia.org/wiki/Canon_de_75_CA_modèle_1940_Schneider}{75~mm Scheider}, \href{https://de.wikipedia.org/wiki/Fliegerabwehrtruppen}{Swiss Fliegerabwehrtruppen}.
            \item militaerfahrzeuge.ch on the \href{https://militaerfahrzeuge.ch/unterkategorie_26_34_333.html}{Flab Kan 38}
        \end{notesources}
    }

    \subsectionwithcounter{KS-12 battery}{KS-12 Battery}

    The KS-12 (M1939 52-K) and the very similar KS-12A (M1944 KS-1) have a single 85~mm gun in an open mount on a towed carriage. A battery has six gun mounts.
    The rate of fire of the KS-12 is 10 rounds per minute.
    Its fire-control system is optical, but the guns are often paired with the SON-9 (Fire Can) radar.
    The KS-12 entered service with Soviet forces in 1939. After WW2, it began to be replaced in Soviet service by the KS-19 100~mm and KS-30 130~mm guns, but was widely used by Soviet allies and saw combat with communist forces in the Korean War and the Vietnam War.%
    \note{KS-12 Battery}{

        \aaamountvalues{85}{10}{790}

        \aaavalues{27}{6/12/18}{2/1/0}{6}{}


        The AAA values are from {\AirStr} and {\TSOH}.

        There are two KS-12 guns. The M-1939 is given a maximum altitude of 34k feet by Wikipedia and 25k feet by Singh. The M-1944 is given a maximum altitude of 38k ft by Wikipedia 29k ft by Singh. The game version have maximum altitudes of 27k feet, which is half-way between the figures given by Singh.

        Wikipedia gives maximum horizontal ranges of 15.65 km and 18 km for the M-1939 and M-1944.

        Wikipedia gives the rate of fire as 10--12 RPM.

        It is not clear to me that the Soviets had proximity fuzes for this gun.

        % LTeX: dictionary+=Defence
        \begin{notesources}
            \item Wikipedia on the \href{https://en.wikipedia.org/wiki/85_mm_air_defense_gun_M1939_(52-K)}{KS-12}.
            \item WeaponSystems on the \href{https://weaponsystems.net/system/1439-KS-12}{KS-12}.
            \item Singh, M. (Pen \& Sword, 2020): “Anti-Aircraft Artillery in Combat 1950--1972: Air Defence in the Jet Age”.
        \end{notesources}
        % LTeX: dictionary-=Defence

    }

    \subsectionwithcounter{90 mm M2 battery}{90 mm M2 Battery}

    The 90 mm M2 has a single 90 mm gun on a towed carriage.
    A battery has four guns.
    The fire control system of the M2 is optical, but an SCR-584 FCR-A normally controls each battery.
    The mechanical loading system permits a rate of fire of 24 rounds per minute.
    Proximity-fuzed shells are available.
    The M2 entered service with the US Army in 1943, following on from the similar M1 in 1940, and served as the main US heavy AA gun in WW2 and the Korean War. It was retired in the mid 1950s. It also served with Canadian (until 1960), French, and Turkish forces.%
    \note{90 mm M2}{

        \aaamountvalues{90}{24}{820}

        \aaavalues{43}{7/14/21}{2/2/1}{6}{}


        The M1 version entered service in 1940. The M2 was a development of the M1 and was the “standard weapon” from 1943. I presume all post-war examples are M2s.

        Wikipedia gives the rate of fire as 24 RPM, about twice that of the KS-12.

        Wikipedia gives a horizontal range of 19 km, and vertical range of 43.5k feet. Werrell states the effective ceiling is 32k feet. As the caliber and the muzzle velocity of the M2 are between the KS-12 (maximum altitude of 27) and the KS-19 (maximum altitude of 39), this seems to be about right.

        I give the M2 the hit rolls of the KS-12 increased by 1, to account for the increased rate of fire and proximity fuzes (although the number of guns per battery is reduced from six to four), a maximum altitude of 32 (following Werrell), and the ranges of the KS-19.

        \begin{notesources}
            \item Wikipedia on the \href{https://en.wikipedia.org/wiki/90_mm_gun_M1/M2/M3}{90 mm M2}.
            \item Werrell, K.P. (2005, Air University Press, 2nd edition): “\href{https://media.defense.gov/2017/Mar/31/2001725225/-1/-1/0/B_0028_WERRELL_ARCHIE_TO_SAM.PDF}{Archie to SAM}”.
        \end{notesources}
    }

    \subsectionwithcounter{QF 3.7-inch battery}{QF 3.7-inch Battery}

    The QF 3.7-inch has a single 94 mm gun on a towed carriage, but typically is used from a prepared firing site.
    The fire control system is optical, but an SCR-584 FCR-A normally controls each battery.
    In the final Mk VI version, the mechanical loading system permits a rate of fire of 20 rounds per minute.
    Proximity-fuzed shells are available.
    The QF 3.7-inch entered service with the British Army in 1937, served as its principal heavy AA gun during WW2 and beyond, and was withdrawn in 1957. It also served with Australia, Belgium, Canada, Sri Lanka, Cyprus, India, Ireland, Israel, Malta, Nepal, New Zealand, Pakistan, South Africa, Yugoslavia, and Portugal.%
    \note{QF 3.7-inch Battery}{

        \aaamountvalues{94}{20}{1044}

        \aaavalues{45}{9/18/27}{3/2/1}{6}{}


        Wikipedia gives a maximum altitude of 45k feet, a muzzle velocity of 1044 m/s (significantly more than other similar guns), and a rate of fire of 20 rounds per minute.

        I give the QF 3.7 the hit rolls of the KS-12 increased by 1, to account for the increased rate of fire and proximity fuzes, a maximum altitude of 45, and a maximum range of 27.

        Singh states that by 1956 the IDF had twelve 3.7-inch guns with FCR.

        Singh states that Pakistan used the 3.7-inch gun in the 1971 War.

        % LTeX: dictionary+=Defence
        \begin{notesources}
            \item Wikipedia on the \href{https://en.wikipedia.org/wiki/QF_3.7-inch_AA_gun}{QF 3.7-inch}.
            \item Singh, M. (Pen \& Sword, 2020): “Anti-Aircraft Artillery in Combat 1950--1972: Air Defence in the Jet Age”.
        \end{notesources}
        % LTeX: dictionary-=Defence

    }

    \subsectionwithcounter{KS-19 battery}{KS-19 Battery}

    The KS-19 has a single 100~mm gun in an open mount on a towed carriage.
    A battery has six gun mounts.
    The fire-control system of the KS-19 is optical, but like other contemporary Soviet guns, it was often paired with the SON-9 (Fire Can) radar.
    The KS-19 entered service in 1948 with Soviet forces as a replacement for the KS-12 85~mm gun. It also served with many Soviet allies and saw combat with communist forces in the Korean War and Vietnam War and Iraqi forces in the Iran-Iraq and Gulf Wars.%
    \note{KS-19 Battery}{

        \aaamountvalues{100}{15}{900}

        \aaavalues{39}{7/14/21}{1/1/0}{6}{}


        The AAA values are from {\AirStr}.

        Wikipedia gives a rate of fire of 15 RPM, a vertical range of 48 altitude levels (15 km, with a proximity fuse), and a horizontal range of 39 hexes (21 km). The {\AirStr} values are 39 altitude levels and 21 hexes, so are cut back somewhat from the maximums.

        It is not clear to me that the Soviets had proximity fuzes for this gun. The one suggestion that they do is that Wikipedia mentions different maximum altitudes for proximity and timed fuzes, but I can find no information on Soviet proximity fuzes.

        \begin{notesources}
            \item Wikipedia on the \href{https://en.wikipedia.org/wiki/KS-19}{KS-19}.
            \item WeaponSystems on the \href{https://weaponsystems.net/system/422-KS-19}{KS-19}.
        \end{notesources}
    }

    \subsectionwithcounter{120 mm M1 battery}{120 mm M1 Battery}

    The M1 has a single 120~mm gun in an open mount on a towed carriage, but typically is used from a prepared firing site. A battery has four guns.
    The fire control system of the M1 is optical, but an SCR-584 FCR-A normally controls each battery. The rate of fire is 12 rounds per minute.
    Proximity-fuzed shells are available.
    The M1 entered service with the US Army in 1944, and though deployed to the Pacific Theater did not engage on enemy aircraft. After the war, it was widely deployed to defend of strategic targets against Soviet Tu-4 bombers, but was withdrawn by 1960. It also served with Canadian forces.%
    \note{120 mm M1 Battery}{

        \aaamountvalues{120}{12}{945}

        \aaavalues{57}{11/22/33}{2/2/1}{7}{}


        Wikipedia gives rate of fire of 12 RPM, a maximum altitude of 57 altitude levels (17.5 km), and a maximum range of 50 hexes (25 km).

        The hit rolls are adapted from the KS-30 (which has a similar rate of fire), with a \plus{1} modifier for proximity shells (although the number of guns per battery is reduced from six to four).

        \begin{notesources}
            \item Wikipedia on the \href{https://en.wikipedia.org/wiki/120_mm_gun_M1}{M1}.
        \end{notesources}
    }

    \subsectionwithcounter{KS-30 battery}{KS-30 Battery}

    The KS-30 has a single 120~mm gun in an open mount on a towed carriage.
    A battery has six gun mounts.
    Development of the KS-30 started after WW2, and it has features adopted from captured German guns as well as British QF 3.7-inch guns.
    The fire-control system is optical, but like other contemporary Soviet guns, it was often paired with the SON-9 (Fire Can) radar. The rate of fire is 12 rounds per minute.
    The KS-30 entered service in 1955 with Soviet forces, was used to defend strategic sites, and served until 1962. It was not widely exported, but also served with the North Vietnamese forces in the Vietnam War and Iraqi forces in the Iran-Iraq and Gulf Wars.%
    \note{KS-30 Battery}{

        \aaamountvalues{130}{12}{970}

        \aaavalues{64}{12/24/36}{1/1/0}{7}{}


        Wikipedia gives the rate of fire is 10--12 RPM, the vertical range as 72 altitude levels (22 km) and 54 hexes (29 km) horizontally. Wikipedia also says it provided a 25\% increase over the effective ceiling of the KS-19, which would be 48 altitude levels. However, this has a heavier shell and slightly higher muzzle velocity than the 120~mm M1, so it should have a maximum altitude of more than 57. I will give it a maximum altitude of 64.

        It is not clear to me that the Soviets had proximity fuzes for this gun.

        The hit rolls are adapted from the KS-19 (which has a similar rate of fire). In order to reach an altitude of 64, it needs a maximum range of at least 32. I will give it 36.

        \begin{notesources}
            \item Wikipedia on the \href{https://en.wikipedia.org/wiki/130_mm_air_defense_gun_KS-30}{KS-30}.
        \end{notesources}
    }

    \subsectionwithcounter{QF 5.25-inch battery}{QF 5.25-inch Battery}

    The QF 5.25-inch has a single 133 mm gun in a fixed, fortified turret.
    It is a dual-purpose gun, used both against aircraft and for coastal defense, adapted from a naval gun.
    A battery has four guns.
    The fire control system is optical, but an SCR-584 or Blue Cedar FCR-A normally controls each battery.
    The hydraulic loading system permits a rate of fire of 10 rounds per minute.
    The QF 5.25-inch entered service with the British Army in 1942 and was installed in and around London. Late in WW2, seven were installed in Australia and three in New Guinea. Eleven were installed in Gibraltar starting in 1953 and were manned until the early 1980s.%
    \note{QF 5.25-inch Battery}{

        \aaamountvalues{133}{10}{850}

        \aaavalues{46}{12/24/36}{2/2/1}{7}{}


        The AAA values are the same as for the KS-30, except the hit rolls are increased by one for proximity fuzes and the maximum altitude is reduced from 64,000 to 46,000 ft. This is the value given on Wikipedia, but it makes sense given the lower muzzle velocity of the 5.25-inch.

        \begin{notesources}
            \item Wikipedia on the \href{https://en.wikipedia.org/wiki/QF_5.25-inch_naval_gun}{QF 5.25-inch} and \href{https://en.wikipedia.org/wiki/Princess_Anne\%27s_Battery}{Princess Anne's Battery}.
        \end{notesources}
    }

\end{withcounters}

\newenvironment{aaaunittable}{
    \onecolumn
    \pagestyle{fancy}
    \raggedbottom
    \footnotesize
    \begin{longtable}{llcccccccccc@{~}c@{~}cc@{~}c@{~}cccccl}

        \longtablefirstcaption{table:aaa-units}{AAA Units Table}\\
        \toprule
        Type&
        \wbox[r]{00}{\vertical{Gun}}&
        \vertical{Country}&
        \vertical{Year}&
        \vertical{Size\tablenotemark{a}}&
        \vertical{Defense Strength}&
        \vertical{Sighting Range}&
        \vertical{Mobility}&
        \begin{tabular}[b]{@{}c@{}}\vertical{VPs}\\\midrule 3D/2D/D\\
        \end{tabular}&
        \vertical{AAA Class}&
        \multirow[b]{-3}{*}{\vertical{Altitude}}&
        \multicolumn{3}{c}{
            \begin{tabular}[b]{@{}ccc@{}}\multicolumn{3}{@{}c@{}}{\vertical{Range}}\\\midrule\wbox[c]{00}{S}&\wbox[c]{00}{M}&\wbox[c]{00}{L}
        \end{tabular}}&
        \multicolumn{3}{c}{
            \begin{tabular}[b]{@{}ccc@{}}\multicolumn{3}{@{}c@{}}{\vertical{Hit}}\\\midrule\wbox[c]{00}{S}&\wbox[c]{00}{M}&\wbox[c]{00}{L}
        \end{tabular}}&
        \vertical{Damage Rating}&
        \vertical{FCR}&
        \vertical{LRF}&
        \vertical{Night IR Sights}&
        \wbox[r]{00}{\vertical{SAM}}\\
        \midrule
        \addlinespace
        \endfirsthead

        \longtablecaption{AAA Units Table}\\
        \toprule
        Type&
        \wbox[r]{00}{\vertical{Gun}}&
        \vertical{Country}&
        \vertical{Year}&
        \vertical{Size\tablenotemark{a}}&
        \vertical{Defense Strength}&
        \vertical{Sighting Range}&
        \vertical{Mobility}&
        \begin{tabular}[b]{@{}c@{}}\vertical{VPs}\\\midrule 3D/2D/D\\
        \end{tabular}&
        \vertical{AAA Class}&
        \multirow[b]{-3}{*}{\vertical{Altitude}}&
        \multicolumn{3}{c}{
            \begin{tabular}[b]{@{}ccc@{}}\multicolumn{3}{@{}c@{}}{\vertical{Range}}\\\midrule\wbox[c]{00}{S}&\wbox[c]{00}{M}&\wbox[c]{00}{L}
        \end{tabular}}&
        \multicolumn{3}{c}{
            \begin{tabular}[b]{@{}ccc@{}}\multicolumn{3}{@{}c@{}}{\vertical{Hit}}\\\midrule\wbox[c]{00}{S}&\wbox[c]{00}{M}&\wbox[c]{00}{L}
        \end{tabular}}&
        \vertical{Damage Rating}&
        \vertical{FCR}&
        \vertical{LRF}&
        \vertical{Night IR Sights}&
        \wbox[r]{00}{\vertical{SAM}}\\
        \midrule
        \addlinespace
        \endhead

        %\addlinespace
        \bottomrule
        \endfoot
    }{
    \end{longtable}
}
% LTeX: enabled=false
\begin{aaaunittable}
M2 &12.7 mm&US&1933&\wbox[l]{Section}{Platoon}&\wbox[l]{0}{3}&\wbox{00}{12}&U&\wbox{00/00/0}{3/2/1}&L&\wbox{00}{5}&\wbox{00}{2}&\wbox{00}{3}&\wbox{00}{4}&\wbox{00}{3}&\wbox{00}{2}&\wbox{00}{1}&\wbox{0}{1}&---&---&---&---\\
DShK &12.7 mm&SU&1938&\wbox[l]{Section}{Platoon}&\wbox[l]{0}{3}&\wbox{00}{12}&U&\wbox{00/00/0}{3/2/1}&L&\wbox{00}{5}&\wbox{00}{2}&\wbox{00}{3}&\wbox{00}{4}&\wbox{00}{3}&\wbox{00}{2}&\wbox{00}{1}&\wbox{0}{1}&---&---&---&---\\
M16 &\binarymultiply{12.7 mm}{4}&US&1944&\wbox[l]{Section}{Section}&\wbox[l]{0}{\underline{3}!}&\wbox{00}{12}&U&\wbox{00/00/0}{6/3}&L&\wbox{00}{5}&\wbox{00}{2}&\wbox{00}{3}&\wbox{00}{4}&\wbox{00}{4}&\wbox{00}{3}&\wbox{00}{2}&\wbox{0}{2}&---&---&---&---\\
M17 &\binarymultiply{12.7 mm}{4}&US&1944&\wbox[l]{Section}{Section}&\wbox[l]{0}{\underline{3}!}&\wbox{00}{12}&U&\wbox{00/00/0}{6/3}&L&\wbox{00}{5}&\wbox{00}{2}&\wbox{00}{3}&\wbox{00}{4}&\wbox{00}{4}&\wbox{00}{3}&\wbox{00}{2}&\wbox{0}{2}&---&---&---&---\\
M55 &\binarymultiply{12.7 mm}{4}&US&1945&\wbox[l]{Section}{Platoon}&\wbox[l]{0}{3}&\wbox{00}{12}&T&\wbox{00/00/0}{4/3/1}&L&\wbox{00}{5}&\wbox{00}{2}&\wbox{00}{3}&\wbox{00}{4}&\wbox{00}{5}&\wbox{00}{4}&\wbox{00}{3}&\wbox{0}{2}&---&---&---&---\\
Vz.53 &\binarymultiply{12.7 mm}{4}&CS&1953&\wbox[l]{Section}{Platoon}&\wbox[l]{0}{3}&\wbox{00}{12}&T&\wbox{00/00/0}{4/3/1}&L&\wbox{00}{5}&\wbox{00}{2}&\wbox{00}{3}&\wbox{00}{4}&\wbox{00}{5}&\wbox{00}{4}&\wbox{00}{3}&\wbox{0}{2}&---&---&---&---\\
BTR-152A Vz.53 &\binarymultiply{12.7 mm}{4}&CS&1955&\wbox[l]{Section}{Section}&\wbox[l]{0}{\underline{3}!}&\wbox{00}{12}&G/P&\wbox{00/00/0}{6/3}&L&\wbox{00}{5}&\wbox{00}{2}&\wbox{00}{3}&\wbox{00}{4}&\wbox{00}{4}&\wbox{00}{3}&\wbox{00}{2}&\wbox{0}{2}&---&---&---&---\\
Mobile M45 &\binarymultiply{12.7 mm}{4}&US&1967&\wbox[l]{Section}{Section}&\wbox[l]{0}{2}&\wbox{00}{12}&G/P&\wbox{00/00/0}{4/2}&L&\wbox{00}{5}&\wbox{00}{2}&\wbox{00}{3}&\wbox{00}{4}&\wbox{00}{4}&\wbox{00}{3}&\wbox{00}{2}&\wbox{0}{2}&---&---&---&---\\
\addlinespace
ZPU-1 &14.5 mm&SU&1949&\wbox[l]{Section}{Battery}&\wbox[l]{0}{3}&\wbox{00}{12}&T&\wbox{00/00/0}{3/2/1}&L&\wbox{00}{6}&\wbox{00}{2}&\wbox{00}{3}&\wbox{00}{5}&\wbox{00}{3}&\wbox{00}{2}&\wbox{00}{1}&\wbox{0}{1}&---&---&---&---\\
ZPU-2 &\binarymultiply{14.5 mm}{2}&SU&1949&\wbox[l]{Section}{Battery}&\wbox[l]{0}{3}&\wbox{00}{12}&T&\wbox{00/00/0}{4/3/1}&L&\wbox{00}{6}&\wbox{00}{2}&\wbox{00}{3}&\wbox{00}{5}&\wbox{00}{4}&\wbox{00}{3}&\wbox{00}{2}&\wbox{0}{1}&---&---&---&---\\
ZPU-4 &\binarymultiply{14.5 mm}{4}&SU&1949&\wbox[l]{Section}{Battery}&\wbox[l]{0}{3}&\wbox{00}{12}&T&\wbox{00/00/0}{5/3/2}&L&\wbox{00}{6}&\wbox{00}{2}&\wbox{00}{3}&\wbox{00}{5}&\wbox{00}{5}&\wbox{00}{4}&\wbox{00}{3}&\wbox{0}{2}&---&---&---&---\\
Mobile ZPU-1 &14.5 mm&SU&1949&\wbox[l]{Section}{Section}&\wbox[l]{0}{2}&\wbox{00}{12}&G/P&\wbox{00/00/0}{2/1}&L&\wbox{00}{6}&\wbox{00}{2}&\wbox{00}{3}&\wbox{00}{5}&\wbox{00}{2}&\wbox{00}{1}&\wbox{00}{0}&\wbox{0}{1}&---&---&---&---\\
Mobile ZPU-2 &\binarymultiply{14.5 mm}{2}&SU&1949&\wbox[l]{Section}{Section}&\wbox[l]{0}{2}&\wbox{00}{12}&G/P&\wbox{00/00/0}{3/2}&L&\wbox{00}{6}&\wbox{00}{2}&\wbox{00}{3}&\wbox{00}{5}&\wbox{00}{3}&\wbox{00}{2}&\wbox{00}{1}&\wbox{0}{1}&---&---&---&---\\
Mobile ZPU-4 &\binarymultiply{14.5 mm}{4}&SU&1949&\wbox[l]{Section}{Section}&\wbox[l]{0}{2}&\wbox{00}{12}&G/P&\wbox{00/00/0}{4/2}&L&\wbox{00}{6}&\wbox{00}{2}&\wbox{00}{3}&\wbox{00}{5}&\wbox{00}{4}&\wbox{00}{3}&\wbox{00}{2}&\wbox{0}{2}&---&---&---&---\\
BTR-40A &\binarymultiply{14.5 mm}{2}&SU&1950&\wbox[l]{Section}{Section}&\wbox[l]{0}{\underline{3}!}&\wbox{00}{12}&G/P&\wbox{00/00/0}{4/2}&L&\wbox{00}{6}&\wbox{00}{2}&\wbox{00}{3}&\wbox{00}{5}&\wbox{00}{3}&\wbox{00}{2}&\wbox{00}{1}&\wbox{0}{1}&---&---&---&---\\
BTR-152A ZPU-2 &\binarymultiply{14.5 mm}{2}&SU&1951&\wbox[l]{Section}{Section}&\wbox[l]{0}{\underline{3}!}&\wbox{00}{12}&G/P&\wbox{00/00/0}{4/2}&L&\wbox{00}{6}&\wbox{00}{2}&\wbox{00}{3}&\wbox{00}{5}&\wbox{00}{3}&\wbox{00}{2}&\wbox{00}{1}&\wbox{0}{1}&---&---&---&---\\
BTR-152A ZPU-4 &\binarymultiply{14.5 mm}{4}&SU&1951&\wbox[l]{Section}{Section}&\wbox[l]{0}{\underline{3}!}&\wbox{00}{12}&G/P&\wbox{00/00/0}{5/3}&L&\wbox{00}{6}&\wbox{00}{2}&\wbox{00}{3}&\wbox{00}{5}&\wbox{00}{4}&\wbox{00}{3}&\wbox{00}{2}&\wbox{0}{2}&---&---&---&---\\
M113 ZPU-2 &\binarymultiply{14.5 mm}{2}&LB&1975&\wbox[l]{Section}{Section}&\wbox[l]{0}{\underline{4}!}&\wbox{00}{12}&U&\wbox{00/00/0}{4/2}&L&\wbox{00}{6}&\wbox{00}{2}&\wbox{00}{3}&\wbox{00}{5}&\wbox{00}{3}&\wbox{00}{2}&\wbox{00}{1}&\wbox{0}{1}&---&---&---&---\\
M113 ZPU-4 &\binarymultiply{14.5 mm}{4}&LB&1975&\wbox[l]{Section}{Section}&\wbox[l]{0}{\underline{4}!}&\wbox{00}{12}&U&\wbox{00/00/0}{5/3}&L&\wbox{00}{6}&\wbox{00}{2}&\wbox{00}{3}&\wbox{00}{5}&\wbox{00}{4}&\wbox{00}{3}&\wbox{00}{2}&\wbox{0}{2}&---&---&---&---\\
\addlinespace
53 T1 &20 mm&FR&1953&\wbox[l]{Section}{Platoon}&\wbox[l]{0}{3}&\wbox{00}{12}&T&\wbox{00/00/0}{4/3/1}&L&\wbox{00}{8}&\wbox{00}{2}&\wbox{00}{4}&\wbox{00}{6}&\wbox{00}{2}&\wbox{00}{2}&\wbox{00}{1}&\wbox{0}{2}&---&---&---&---\\
M167 &\binarymultiply{20 mm}{6}&US&1967&\wbox[l]{Section}{Platoon}&\wbox[l]{0}{3}&\wbox{00}{12}&T&\wbox{00/00/0}{7/5/2}&L&\wbox{00}{8}&\wbox{00}{2}&\wbox{00}{4}&\wbox{00}{6}&\wbox{00}{6}&\wbox{00}{5}&\wbox{00}{4}&\wbox{0}{3}&R/HF&---&---&---\\
M163 &\binarymultiply{20 mm}{6}&US&1968&\wbox[l]{Section}{Section}&\wbox[l]{0}{\underline{4}!}&\wbox{00}{12}&U&\wbox{00/00/0}{8/4}&L&\wbox{00}{8}&\wbox{00}{2}&\wbox{00}{4}&\wbox{00}{6}&\wbox{00}{5}&\wbox{00}{4}&\wbox{00}{3}&\wbox{0}{3}&R/HF&---&---&---\\
TCM-20 &\binarymultiply{20 mm}{2}&IL&1970&\wbox[l]{Section}{Platoon}&\wbox[l]{0}{3}&\wbox{00}{12}&T&\wbox{00/00/0}{5/3/2}&L&\wbox{00}{8}&\wbox{00}{2}&\wbox{00}{4}&\wbox{00}{6}&\wbox{00}{4}&\wbox{00}{4}&\wbox{00}{3}&\wbox{0}{2}&---&---&---&---\\
M3 TCM-20 &\binarymultiply{20 mm}{2}&IL&1970&\wbox[l]{Section}{Section}&\wbox[l]{0}{\underline{3}!}&\wbox{00}{12}&U&\wbox{00/00/0}{6/3}&L&\wbox{00}{8}&\wbox{00}{2}&\wbox{00}{4}&\wbox{00}{6}&\wbox{00}{3}&\wbox{00}{3}&\wbox{00}{2}&\wbox{0}{2}&---&---&---&---\\
Rh-202 &\binarymultiply{20 mm}{2}&DE&1972&\wbox[l]{Section}{Battery}&\wbox[l]{0}{3}&\wbox{00}{12}&T&\wbox{00/00/0}{5/3/2}&L&\wbox{00}{8}&\wbox{00}{2}&\wbox{00}{4}&\wbox{00}{6}&\wbox{00}{5}&\wbox{00}{4}&\wbox{00}{3}&\wbox{0}{3}&---&---&---&---\\
Panhard M3 VDA &\binarymultiply{20 mm}{2}&FR&1975&\wbox[l]{Section}{Section}&\wbox[l]{0}{\underline{4}}&\wbox{00}{12}&U&\wbox{00/00/0}{7/4}&L&\wbox{00}{8}&\wbox{00}{2}&\wbox{00}{4}&\wbox{00}{6}&\wbox{00}{4}&\wbox{00}{4}&\wbox{00}{3}&\wbox{0}{2}&R/UF&---&---&---\\
53 T1 M693 &20 mm&FR&1976&\wbox[l]{Section}{Platoon}&\wbox[l]{0}{3}&\wbox{00}{12}&T&\wbox{00/00/0}{4/3/1}&L&\wbox{00}{8}&\wbox{00}{2}&\wbox{00}{4}&\wbox{00}{6}&\wbox{00}{3}&\wbox{00}{3}&\wbox{00}{2}&\wbox{0}{2}&---&---&---&---\\
Tarasque &20 mm&FR&1981&\wbox[l]{Section}{Platoon}&\wbox[l]{0}{3}&\wbox{00}{12}&T&\wbox{00/00/0}{4/3/1}&L&\wbox{00}{8}&\wbox{00}{2}&\wbox{00}{4}&\wbox{00}{6}&\wbox{00}{4}&\wbox{00}{3}&\wbox{00}{2}&\wbox{0}{2}&---&---&---&---\\
Mobile Tarasque &20 mm&FR&1982&\wbox[l]{Section}{Section}&\wbox[l]{0}{2}&\wbox{00}{12}&G/P&\wbox{00/00/0}{3/2}&L&\wbox{00}{8}&\wbox{00}{2}&\wbox{00}{4}&\wbox{00}{6}&\wbox{00}{3}&\wbox{00}{2}&\wbox{00}{1}&\wbox{0}{2}&---&---&---&---\\
\addlinespace
ZU-23 &\binarymultiply{23 mm}{2}&SU&1960&\wbox[l]{Section}{Battery}&\wbox[l]{0}{3}&\wbox{00}{12}&T&\wbox{00/00/0}{6/4/2}&L&\wbox{00}{9}&\wbox{00}{2}&\wbox{00}{5}&\wbox{00}{8}&\wbox{00}{5}&\wbox{00}{4}&\wbox{00}{3}&\wbox{0}{3}&---&---&---&---\\
Mobile ZU-23 &\binarymultiply{23 mm}{2}&SU&1960&\wbox[l]{Section}{Section}&\wbox[l]{0}{2}&\wbox{00}{12}&G/P&\wbox{00/00/0}{5/3}&L&\wbox{00}{9}&\wbox{00}{2}&\wbox{00}{5}&\wbox{00}{8}&\wbox{00}{4}&\wbox{00}{3}&\wbox{00}{2}&\wbox{0}{3}&---&---&---&---\\
ZSU-23-4 &\binarymultiply{23 mm}{4}&SU&1965&\wbox[l]{Section}{Section}&\wbox[l]{0}{\underline{4}}&\wbox{00}{12}&U&\wbox{00/00/0}{9/5}&L&\wbox{00}{9}&\wbox{00}{2}&\wbox{00}{5}&\wbox{00}{8}&\wbox{00}{5}&\wbox{00}{4}&\wbox{00}{3}&\wbox{0}{3}&W/VF&---&---&---\\
BTR-152A ZU-23 &\binarymultiply{23 mm}{2}&PS&1975&\wbox[l]{Section}{Section}&\wbox[l]{0}{\underline{3}!}&\wbox{00}{12}&G/P&\wbox{00/00/0}{4/2}&L&\wbox{00}{9}&\wbox{00}{2}&\wbox{00}{5}&\wbox{00}{8}&\wbox{00}{4}&\wbox{00}{3}&\wbox{00}{2}&\wbox{0}{3}&---&---&---&---\\
Sinai 23 Ayn-al-Saqr &\binarymultiply{23 mm}{2}&EG&1992&\wbox[l]{Section}{Section}&\wbox[l]{0}{\underline{4}}&\wbox{00}{12}&U&\wbox{00/00/0}{10/5}&L&\wbox{00}{9}&\wbox{00}{2}&\wbox{00}{5}&\wbox{00}{8}&\wbox{00}{5}&\wbox{00}{4}&\wbox{00}{3}&\wbox{0}{3}&R/UF&---&---&Ayn-al-Saqr\\
Sinai 23 Stinger &\binarymultiply{23 mm}{2}&EG&1992&\wbox[l]{Section}{Section}&\wbox[l]{0}{\underline{4}}&\wbox{00}{12}&U&\wbox{00/00/0}{12/6}&L&\wbox{00}{9}&\wbox{00}{2}&\wbox{00}{5}&\wbox{00}{8}&\wbox{00}{5}&\wbox{00}{4}&\wbox{00}{3}&\wbox{0}{3}&R/UF&---&---&FIM-92\\
\addlinespace
M2 Blazer &\binarymultiply{25 mm}{5}&US&---&\wbox[l]{Section}{Section}&\wbox[l]{0}{\underline{4}}&\wbox{00}{12}&U&\wbox{00/00/0}{14/7}&L&\wbox{00}{10}&\wbox{00}{3}&\wbox{00}{6}&\wbox{00}{9}&\wbox{00}{5}&\wbox{00}{4}&\wbox{00}{3}&\wbox{0}{4}&S/VF&Y&Y&FIM-92D\\
LAV-AD &\binarymultiply{25 mm}{5}&US&1997&\wbox[l]{Section}{Section}&\wbox[l]{0}{\underline{4}}&\wbox{00}{12}&U&\wbox{00/00/0}{14/7}&L&\wbox{00}{10}&\wbox{00}{3}&\wbox{00}{6}&\wbox{00}{9}&\wbox{00}{5}&\wbox{00}{4}&\wbox{00}{3}&\wbox{0}{4}&---&Y&Y&FIM-92D\\
\addlinespace
M53/59 &\binarymultiply{30 mm}{2}&CS&1959&\wbox[l]{Section}{Section}&\wbox[l]{0}{\underline{4}!}&\wbox{00}{12}&G/P&\wbox{00/00/0}{7/4}&L&\wbox{00}{11}&\wbox{00}{4}&\wbox{00}{7}&\wbox{00}{9}&\wbox{00}{3}&\wbox{00}{2}&\wbox{00}{1}&\wbox{0}{4}&---&---&---&---\\
HS.661A &30 mm&CH&1955&\wbox[l]{Section}{Battery}&\wbox[l]{0}{3}&\wbox{00}{12}&T&\wbox{00/00/0}{8/5/3}&L&\wbox{00}{11}&\wbox{00}{4}&\wbox{00}{7}&\wbox{00}{9}&\wbox{00}{2}&\wbox{00}{2}&\wbox{00}{1}&\wbox{0}{3}&---&---&---&---\\
HS.661B &30 mm&CH&1960&\wbox[l]{Section}{Battery}&\wbox[l]{0}{3}&\wbox{00}{12}&T&\wbox{00/00/0}{8/5/3}&L&\wbox{00}{11}&\wbox{00}{4}&\wbox{00}{7}&\wbox{00}{9}&\wbox{00}{3}&\wbox{00}{2}&\wbox{00}{1}&\wbox{0}{3}&---&---&---&---\\
AMX-13 DCA &\binarymultiply{30 mm}{2}&FR&1969&\wbox[l]{Section}{Section}&\wbox[l]{0}{\underline{4}}&\wbox{00}{12}&U&\wbox{00/00/0}{10/5}&L&\wbox{00}{11}&\wbox{00}{4}&\wbox{00}{7}&\wbox{00}{9}&\wbox{00}{4}&\wbox{00}{3}&\wbox{00}{2}&\wbox{0}{3}&R/VF&---&---&---\\
AMX-30 DCA &\binarymultiply{30 mm}{2}&FR&---&\wbox[l]{Section}{Section}&\wbox[l]{0}{\underline{4}}&\wbox{00}{12}&U&\wbox{00/00/0}{10/5}&L&\wbox{00}{11}&\wbox{00}{4}&\wbox{00}{7}&\wbox{00}{9}&\wbox{00}{4}&\wbox{00}{3}&\wbox{00}{2}&\wbox{0}{3}&R/VF&---&---&---\\
AMX-30SA &\binarymultiply{30 mm}{2}&FR&1979&\wbox[l]{Section}{Section}&\wbox[l]{0}{\underline{4}}&\wbox{00}{12}&U&\wbox{00/00/0}{10/5}&L&\wbox{00}{11}&\wbox{00}{4}&\wbox{00}{7}&\wbox{00}{9}&\wbox{00}{5}&\wbox{00}{4}&\wbox{00}{3}&\wbox{0}{3}&R/VF&---&---&---\\
Tunguska &\binarymultiply{30 mm}{2}&SU&1984&\wbox[l]{Section}{Section}&\wbox[l]{0}{\underline{4}}&\wbox{00}{12}&U&\wbox{00/00/0}{20/10}&L&\wbox{00}{11}&\wbox{00}{4}&\wbox{00}{7}&\wbox{00}{9}&\wbox{00}{5}&\wbox{00}{4}&\wbox{00}{3}&\wbox{0}{4}&W/VF&---&---&SA-19\\
Tunguska-M &\binarymultiply{30 mm}{2}&SU&1990&\wbox[l]{Section}{Section}&\wbox[l]{0}{\underline{4}}&\wbox{00}{12}&U&\wbox{00/00/0}{22/11}&L&\wbox{00}{11}&\wbox{00}{4}&\wbox{00}{7}&\wbox{00}{9}&\wbox{00}{5}&\wbox{00}{4}&\wbox{00}{3}&\wbox{0}{4}&W/VF&---&---&SA-19\\
\pagebreak
Oerlikon GDF-001 &\binarymultiply{35 mm}{2}&CH&1963&\wbox[l]{Section}{Battery}&\wbox[l]{0}{3}&\wbox{00}{12}&T&\wbox{00/00/0}{7/5/2}&M&\wbox{00}{15}&\wbox{00}{4}&\wbox{00}{8}&\wbox{00}{12}&\wbox{00}{4}&\wbox{00}{3}&\wbox{00}{2}&\wbox{0}{4}&---&---&---&---\\
Gepard &\binarymultiply{35 mm}{2}&DE&1976&\wbox[l]{Section}{Section}&\wbox[l]{0}{\underline{4}}&\wbox{00}{12}&U&\wbox{00/00/0}{11/6}&M&\wbox{00}{15}&\wbox{00}{4}&\wbox{00}{8}&\wbox{00}{12}&\wbox{00}{5}&\wbox{00}{4}&\wbox{00}{3}&\wbox{0}{4}&W/VF&---&---&---\\
Oerlikon GDF-005 &\binarymultiply{35 mm}{2}&CH&1985&\wbox[l]{Section}{Battery}&\wbox[l]{0}{3}&\wbox{00}{12}&T&\wbox{00/00/0}{8/5/3}&M&\wbox{00}{15}&\wbox{00}{4}&\wbox{00}{8}&\wbox{00}{12}&\wbox{00}{5}&\wbox{00}{4}&\wbox{00}{3}&\wbox{0}{4}&---&Y&---&---\\
\addlinespace
61-K &37 mm&SU&1939&\wbox[l]{Section}{Battery}&\wbox[l]{0}{3}&\wbox{00}{12}&T&\wbox{00/00/0}{5/3/2}&M&\wbox{00}{12}&\wbox{00}{4}&\wbox{00}{8}&\wbox{00}{11}&\wbox{00}{3}&\wbox{00}{2}&\wbox{00}{1}&\wbox{0}{4}&---&---&---&---\\
M15 &\binaryplus{37 mm}{\binarymultiply{12.7 mm}{2}}&US&1942&\wbox[l]{Section}{Section}&\wbox[l]{0}{\underline{3}!}&\wbox{00}{12}&U&\wbox{00/00/0}{4/2}&M&\wbox{00}{12}&\wbox{00}{2}&\wbox{00}{4}&\wbox{00}{8}&\wbox{00}{3}&\wbox{00}{2}&\wbox{00}{1}&\wbox{0}{4}&---&---&---&---\\
\addlinespace
Bofors L/60 &40 mm&SE&1935&\wbox[l]{Section}{Battery}&\wbox[l]{0}{3}&\wbox{00}{12}&T&\wbox{00/00/0}{5/3/2}&M&\wbox{00}{15}&\wbox{00}{4}&\wbox{00}{8}&\wbox{00}{12}&\wbox{00}{2}&\wbox{00}{2}&\wbox{00}{1}&\wbox{0}{4}&---&---&---&---\\
M19 &\binarymultiply{40 mm}{2}&US&1944&\wbox[l]{Section}{Section}&\wbox[l]{0}{\underline{3}!}&\wbox{00}{12}&U&\wbox{00/00/0}{6/3}&M&\wbox{00}{15}&\wbox{00}{4}&\wbox{00}{8}&\wbox{00}{12}&\wbox{00}{2}&\wbox{00}{2}&\wbox{00}{1}&\wbox{0}{4}&---&---&---&---\\
M34 &40 mm&US&1951&\wbox[l]{Section}{Section}&\wbox[l]{0}{2}&\wbox{00}{12}&U&\wbox{00/00/0}{3/2}&M&\wbox{00}{15}&\wbox{00}{2}&\wbox{00}{4}&\wbox{00}{8}&\wbox{00}{1}&\wbox{00}{1}&\wbox{00}{1}&\wbox{0}{4}&---&---&---&---\\
Bofors L/70 &40 mm&SE&1952&\wbox[l]{Section}{Battery}&\wbox[l]{0}{3}&\wbox{00}{12}&T&\wbox{00/00/0}{6/4/2}&M&\wbox{00}{15}&\wbox{00}{4}&\wbox{00}{8}&\wbox{00}{12}&\wbox{00}{3}&\wbox{00}{2}&\wbox{00}{1}&\wbox{0}{4}&---&---&---&---\\
M42 &\binarymultiply{40 mm}{2}&US&1953&\wbox[l]{Section}{Section}&\wbox[l]{0}{\underline{3}!}&\wbox{00}{12}&U&\wbox{00/00/0}{6/3}&M&\wbox{00}{15}&\wbox{00}{4}&\wbox{00}{8}&\wbox{00}{12}&\wbox{00}{2}&\wbox{00}{2}&\wbox{00}{1}&\wbox{0}{4}&---&---&---&---\\
Bofors L/70 BOFI &40 mm&SE&---&\wbox[l]{Section}{Battery}&\wbox[l]{0}{3}&\wbox{00}{12}&T&\wbox{00/00/0}{8/5/3}&M&\wbox{00}{15}&\wbox{00}{4}&\wbox{00}{8}&\wbox{00}{12}&\wbox{00}{5}&\wbox{00}{4}&\wbox{00}{3}&\wbox{0}{4}&---&---&---&---\\
Bofors L/70 BOFI-R &40 mm&SE&---&\wbox[l]{Section}{Battery}&\wbox[l]{0}{3}&\wbox{00}{12}&T&\wbox{00/00/0}{9/6/3}&M&\wbox{00}{15}&\wbox{00}{4}&\wbox{00}{8}&\wbox{00}{12}&\wbox{00}{6}&\wbox{00}{5}&\wbox{00}{4}&\wbox{0}{4}&W/VF&---&---&---\\
\addlinespace
S-60 &57 mm&SU&1950&\wbox[l]{Section}{Battery}&\wbox[l]{0}{3}&\wbox{00}{12}&T&\wbox{00/00/0}{8/5/3}&M&\wbox{00}{18}&\wbox{00}{5}&\wbox{00}{10}&\wbox{00}{15}&\wbox{00}{2}&\wbox{00}{2}&\wbox{00}{1}&\wbox{0}{5}&---&---&---&---\\
ZSU-57-2 &\binarymultiply{57 mm}{2}&SU&1955&\wbox[l]{Section}{Section}&\wbox[l]{0}{\underline{4}!}&\wbox{00}{12}&U&\wbox{00/00/0}{6/3}&M&\wbox{00}{18}&\wbox{00}{5}&\wbox{00}{10}&\wbox{00}{15}&\wbox{00}{2}&\wbox{00}{1}&\wbox{00}{1}&\wbox{0}{5}&---&---&---&---\\
\addlinespace
Flab Kan 38 &75 mm&CH&1938&\wbox[l]{Section}{Battery}&\wbox[l]{0}{4}&\wbox{00}{18}&T&\wbox{00/00/0}{9/6/3}&H&\wbox{00}{24}&\wbox{00}{6}&\wbox{00}{12}&\wbox{00}{18}&\wbox{00}{1}&\wbox{00}{1}&\wbox{00}{0}&\wbox{0}{5}&---&---&---&---\\
M51 &75 mm&US&1953&\wbox[l]{Section}{Battery}&\wbox[l]{0}{3}&\wbox{00}{18}&T&\wbox{00/00/0}{9/6/3}&M&\wbox{00}{30}&\wbox{00}{6}&\wbox{00}{12}&\wbox{00}{18}&\wbox{00}{4}&\wbox{00}{4}&\wbox{00}{3}&\wbox{0}{5}&W/MF&---&---&---\\
\addlinespace
KS-12 &85 mm&SU&1939&\wbox[l]{Section}{Battery}&\wbox[l]{0}{4}&\wbox{00}{18}&T&\wbox{00/00/0}{9/6/3}&H&\wbox{00}{27}&\wbox{00}{6}&\wbox{00}{12}&\wbox{00}{18}&\wbox{00}{1}&\wbox{00}{1}&\wbox{00}{0}&\wbox{0}{6}&---&---&---&---\\
\addlinespace
90 mm M2 &90 mm&US&1943&\wbox[l]{Section}{Battery}&\wbox[l]{0}{4}&\wbox{00}{18}&T&\wbox{00/00/0}{9/6/3}&H&\wbox{00}{43}&\wbox{00}{7}&\wbox{00}{14}&\wbox{00}{21}&\wbox{00}{2}&\wbox{00}{2}&\wbox{00}{1}&\wbox{0}{6}&---&---&---&---\\
\addlinespace
QF 3.7-Inch &94 mm&UK&1939&\wbox[l]{Section}{Battery}&\wbox[l]{0}{4}&\wbox{00}{18}&T&\wbox{00/00/0}{9/6/3}&H&\wbox{00}{45}&\wbox{00}{9}&\wbox{00}{18}&\wbox{00}{27}&\wbox{00}{2}&\wbox{00}{2}&\wbox{00}{1}&\wbox{0}{6}&---&---&---&---\\
\addlinespace
KS-19 &100 mm&SU&1948&\wbox[l]{Section}{Battery}&\wbox[l]{0}{4}&\wbox{00}{18}&T&\wbox{00/00/0}{10/7/3}&H&\wbox{00}{39}&\wbox{00}{7}&\wbox{00}{14}&\wbox{00}{21}&\wbox{00}{1}&\wbox{00}{1}&\wbox{00}{0}&\wbox{0}{6}&---&---&---&---\\
\addlinespace
120 mm M1 &120 mm&US&1944&\wbox[l]{Section}{Battery}&\wbox[l]{0}{4}&\wbox{00}{18}&T&\wbox{00/00/0}{9/6/3}&H&\wbox{00}{57}&\wbox{00}{11}&\wbox{00}{22}&\wbox{00}{33}&\wbox{00}{2}&\wbox{00}{2}&\wbox{00}{1}&\wbox{0}{7}&---&---&---&---\\
\addlinespace
KS-30 &130 mm&SU&1955&\wbox[l]{Section}{Battery}&\wbox[l]{0}{4}&\wbox{00}{18}&T&\wbox{00/00/0}{11/7/4}&H&\wbox{00}{64}&\wbox{00}{12}&\wbox{00}{24}&\wbox{00}{36}&\wbox{00}{1}&\wbox{00}{1}&\wbox{00}{0}&\wbox{0}{7}&---&---&---&---\\
\addlinespace
QF 5.25-Inch &133 mm&UK&1942&\wbox[l]{Section}{Battery}&\wbox[l]{0}{4}&\wbox{00}{18}&T&\wbox{00/00/0}{11/7/4}&H&\wbox{00}{46}&\wbox{00}{12}&\wbox{00}{24}&\wbox{00}{36}&\wbox{00}{2}&\wbox{00}{2}&\wbox{00}{1}&\wbox{0}{7}&---&---&---&---\\
\end{aaaunittable}


\twocolumn
\pagestyle{fancy}
